\documentclass[11pt,a4paper]{article}
\usepackage[T1]{fontenc}
\usepackage[utf8]{inputenc}
\usepackage{lmodern,amsmath,amssymb,amsthm,mathtools,mathrsfs}
\usepackage[margin=26mm,headheight=15pt]{geometry}
\usepackage{microtype,booktabs,longtable,tabularx,array,enumitem}
\usepackage{xcolor,listings,fancyhdr,titlesec,xurl}
\definecolor{inkblue}{RGB}{28,55,83}
\definecolor{quietgray}{RGB}{83,91,99}
\definecolor{papergray}{RGB}{247,248,250}
\usepackage[colorlinks=true,linkcolor=inkblue,citecolor=inkblue,urlcolor=inkblue,
 bookmarksnumbered=true,pdfencoding=auto]{hyperref}
\hypersetup{pdftitle={Euler's Constant and Shifted Mills Constants: Arithmetic Structure, Reciprocal Series, and Proof-Transfer Limits},
 pdfauthor={Keio and Shaitetsu},pdfsubject={Integrated research manuscript covering twenty ranked results and audit items}}
\pagestyle{fancy}
\fancyhf{}
\fancyhead[L]{\small\color{quietgray}Euler's Constant and Shifted Mills Constants}
\fancyhead[R]{\small\color{quietgray}Keio Research Notes}
\fancyfoot[C]{\thepage}
\renewcommand{\headrulewidth}{0.3pt}
\titleformat{\section}{\Large\bfseries\color{inkblue}}{\thesection}{0.7em}{}
\titleformat{\subsection}{\large\bfseries\color{inkblue}}{\thesubsection}{0.65em}{}
\titleformat{\subsubsection}{\normalsize\bfseries}{\thesubsubsection}{0.6em}{}
\setlength{\emergencystretch}{3em}
\setlength{\parskip}{3pt}
\setlist{itemsep=3pt,topsep=5pt}
\setcounter{tocdepth}{1}
\numberwithin{equation}{section}
\newtheorem{theorem}{Theorem}[section]
\newtheorem{lemma}[theorem]{Lemma}
\newtheorem{proposition}[theorem]{Proposition}
\newtheorem{corollary}[theorem]{Corollary}
\theoremstyle{definition}
\newtheorem{definition}[theorem]{Definition}
\newtheorem{hypothesis}[theorem]{Hypothesis}
\theoremstyle{remark}
\newtheorem{remark}[theorem]{Remark}
\newtheorem{example}[theorem]{Example}
\newcommand{\A}{\overline{\mathbb Q}}
\newcommand{\Q}{\mathbb Q}
\newcommand{\Z}{\mathbb Z}
\newcommand{\N}{\mathbb N}
\newcommand{\R}{\mathbb R}
\newcommand{\C}{\mathbb C}
\newcommand{\Er}{\mathcal E}
\newcommand{\K}{\mathcal K}
\newcommand{\Ein}{\operatorname{Ein}}
\newcommand{\Eone}{E_1}
\newcommand{\Tr}{\operatorname{Tr}}
\newcommand{\Norm}{\operatorname{N}}
\newcommand{\vp}{v_p}
\newcommand{\floor}[1]{\left\lfloor #1\right\rfloor}
\newcommand{\ceil}[1]{\left\lceil #1\right\rceil}
\DeclareMathOperator{\Frac}{Frac}
\DeclareMathOperator{\trdeg}{trdeg}
\DeclareMathOperator{\den}{den}
\DeclareMathOperator{\ord}{ord}
\DeclareMathOperator{\Res}{Res}
\lstdefinestyle{certificate}{language=Python,basicstyle=\ttfamily\fontsize{8.5}{10.2}\selectfont,emptylines=0,
 keywordstyle=\color{inkblue}\bfseries,commentstyle=\color{quietgray},
 stringstyle=\color{inkblue},backgroundcolor=\color{papergray},
 frame=single,rulecolor=\color{quietgray!35},framerule=0.3pt,
 breaklines=true,breakatwhitespace=false,columns=fullflexible,
 keepspaces=true,showstringspaces=false,tabsize=4,
 aboveskip=9pt,belowskip=9pt,xleftmargin=5pt,xrightmargin=5pt}
\title{\color{inkblue}\textbf{Euler's Constant and\\Shifted Mills Constants}\\[6pt]
\large Arithmetic Structure, Reciprocal Series,\\and Proof-Transfer Limits}
\author{Keio \and Shaitetsu}
\date{5 November 2026\\[3pt]\small An integrated account of twenty ranked research items}
\begin{document}
\maketitle
\thispagestyle{empty}
\begin{abstract}
We integrate ten gamma-related and ten Mills-related results, reductions,
and audits into one research manuscript.  On the gamma side, winding
versions of the canonical phase decomposition of $E$-function values lead
to relative algebraic-closure statements for fields generated by
exponential-integral values, a radical obstruction involving the
Euler--Gompertz constant, and an explicit failure of relative closure
under field composita.  On the Mills side, shifted exponent clocks give
escape from fixed polynomial relations, while hypothetical algebraicity
of the least classical Mills constant imposes cubic trace identities,
divisible prime gaps, and a transcendental gap exponent.  A common
integer-tail method proves linear independence of two reciprocal series
for each shifted clock and yields a precise irrationality transfer
criterion.  For rational approximations to Euler's constant, we distinguish
positive analytic errors from primitive arithmetic errors, prove a
slow-degree denominator barrier, and record an exact check of 384
higher-degree Pad\'e approximants.  Proofs, explicit hypotheses, provenance,
and reproducible integer-arithmetic certificates are kept together.
The irrationality of Euler's constant and unconditional transcendence of
the classical Mills constant are not established here.
\end{abstract}
\vfill
\noindent\textbf{Keywords.} Euler--Mascheroni constant; Mills constant;
$E$-function values; relative algebraic closure; reciprocal series;
Pisot numbers; Pad\'e approximation; exact arithmetic.
\clearpage
\tableofcontents
\clearpage

\section{Scope, notation, and the twenty-item register}
\label{intro:scope}

This article consolidates the two ten-item rankings developed in the
preceding research dialogue.  The ranking labels G1--G10 and M1--M10 are
retained as a navigation device.  They are not twenty claims of new
literature priority.  The body reorganizes them by mathematical dependency
and supplies proofs or precisely attributed inputs.  Computational checks
are distinguished from infinite statements, and conditional implications
are distinguished from verification of their antecedents.

\subsection{The two original targets}

Euler's constant is
\[
 \gamma=\lim_{n\to\infty}(H_n-\log n),\qquad
 H_n=\sum_{k=1}^n\frac1k.
\]
For a positive exponent sequence $C=(C_k)_{k\ge1}$, define
\[
 \mathcal W(C)=\{A>1:\floor{A^{C_k}}\text{ is prime for every }k\ge1\},
 \qquad \xi(C)=\min\mathcal W(C)
\]
when the minimum exists.  The classical least Mills constant is
$M=\xi((3^k)_{k\ge1})$.  Its irrationality is an external result of
Saito~\cite{Saito2025}.  The conditional transcendence result under the
Density Hypothesis is also external~\cite{SaitoVariants}.

The four shifted constants considered in the dialogue are
\[
 A_\pm=\xi((3^k\pm1)_{k\ge1}),\qquad
 B_\pm=\xi((3^k\pm k)_{k\ge1}).
\]
Their individual transcendence is used as an input from the named working
manuscripts \cite{VMRC064,VMRC008}; those long source proofs are not
reproduced as if they were consequences of Saito's published variants
theorem.  The floor-sequence theorems below state their exact hypotheses
on the base $A$, so their logical dependence is explicit.  Throughout,
$3^k\pm k$ means a \emph{linear} shift, not a fixed integer offset.

\subsection{Analytic notation and a necessary distinction}

Write $\A=\overline{\Q}$, let $\Er$ denote the ring of values of
$E$-functions at algebraic points, and put $\K=\Frac(\Er)$.  An
$E$-function is understood in the standard arithmetic sense used in
\cite{FR2026,Beukers2006}.  The rings $\Er$ and $\K$ must not be
identified: the former is not a field~\cite{FR2026}.
For $a>0$,
\[
 \Ein(a)=\int_0^a\frac{1-e^{-t}}t\,dt,
 \qquad E_1(a)=\int_a^\infty\frac{e^{-t}}t\,dt,
\]
and the standard connection formula~\cite{NISTDLMF} is
\[
 \Ein(a)=\gamma+\log a+E_1(a).
\]
In particular,
\begin{equation}\label{intro:eta-delta}
 \eta:=\Ein(1)=\gamma+\frac\delta e,
 \qquad
 \delta:=\int_0^\infty\frac{e^{-t}}{1+t}\,dt=eE_1(1).
\end{equation}

\clearpage
\subsection{Status vocabulary and provenance}

\emph{Proved here} means that this manuscript contains a proof from its
stated standard inputs; it does not assert publication priority.
\emph{Conditional} means that the implication is established while a
displayed hypothesis is not supplied.  \emph{Finite certificate} means
that a specified finite computation has a reproducible exact verifier.
\emph{Imported} marks a cited theorem or a source-corpus input.
The corpus documents are working manuscripts rather than independent
external validation of the conclusions.  Standard external results and
the authors' source computations are identified separately in the
bibliography.  No theorem below asserts either of the two original
target conclusions unconditionally.

The register on the next pages preserves every item in the user's two
rankings.  The item identifiers in the register link to the corresponding main
results. Formula numbers and theorem labels provide the detailed dependencies.

\subsection{Gamma register: G1--G10}
\label{intro:gamma-register}
\small
\begin{longtable}{@{}p{0.06\textwidth}>{\raggedright\arraybackslash}p{0.41\textwidth}>{\raggedright\arraybackslash}p{0.45\textwidth}@{}}
\toprule
ID & Content retained & Status and precise reach\\
\midrule\endhead
\hyperref[gf:jointclosure]{G1} & Relative algebraic closure of fields generated by positive algebraic-scale $\Ein$ values & Proved using winding and finite-phase descent.  It concerns $\K$, not a decision of the arithmetic nature of $\gamma$.\\[6pt]
\hyperref[gf:delta-roots]{G2} & $\gamma\in\A$ forces $\delta^{1/n}\notin\K$ for $n\ge2$ & Proved implication.  The reverse-use membership hypothesis needed to conclude transcendence of $\gamma$ is not supplied.\\[6pt]
\hyperref[gf:compositum]{G3} & Failure of relative algebraic closure under composita & Explicit counterexample using $\A(e(\eta-q))$ for two distinct $q\in\A$.\\[6pt]
\hyperref[gf:einintersection]{G4} & $\A(\Ein(a))\cap\Er=\A[\Ein(a)]$ & Integral-intersection argument for every nonzero algebraic $a$; distinguishes the value ring from its fraction field.\\[6pt]
\hyperref[gf:epair]{G5} & Algebraic independence of $e,\eta$ and conditional constraints on $\delta$ & Reproved with the winding map.  Conditional complexity of $\delta$ is not a contradiction to algebraicity of $\gamma$.\\[6pt]
\hyperref[rs:transfer]{G6} & Irrationality transfer from negative-clock reciprocal expansions & Proved for bounded integer coefficients with infinitely many nonzero terms.  The requisite identity for $\gamma$ is not known here.\\[6pt]
\hyperref[ga:logfree]{G7} & A log-free positive Pad\'e approximation family & Infinite, unconditional rational family with strictly positive error.  Primitive denominator times error remains the deciding condition.\\[6pt]
\hyperref[ga:slow-barrier]{G8} & Slow-degree primitive denominator barrier & Infinite-range divergence result in its stated growth regime.  It is not a theorem excluding all degrees.\\[6pt]
\hyperref[ga:lattice]{G9} & Exact CMF lattice and normalization costs & A specified lattice and certificate method are analyzed.  Global matrix content is distinguished from the gcd of an individual output pair.\\[6pt]
\hyperref[ga:384]{G10} & 384 higher-degree primitive arithmetic checks & Exact finite certificate, with code included.  No extrapolation to all parameter pairs is made.\\
\bottomrule
\end{longtable}
\normalsize

\clearpage
\subsection{Mills register: M1--M10}
\label{intro:mills-register}
\small
\begin{longtable}{@{}p{0.06\textwidth}>{\raggedright\arraybackslash}p{0.41\textwidth}>{\raggedright\arraybackslash}p{0.45\textwidth}@{}}
\toprule
ID & Content retained & Status and precise reach\\
\midrule\endhead
\hyperref[mi:affine]{M1} & Two-coordinate and three-coordinate escape from fixed algebraic relations & Proved for transcendental bases.  Fixed-polynomial escape is not algebraic independence of individual integer tuples.\\[6pt]
\hyperref[rs:independence]{M2} & Linear independence of $1,\sum p_k^{-1},\sum p_k^{-2}$ & Proved separately for each clock.  The negative clocks need no transcendence assumption; $A^2\notin\Q$ suffices for the positive clocks.\\[6pt]
\hyperref[mi:tail-limit]{M3} & Continuity at the unshifted clock and normalized tail constants & A genuine limiting connection to $M$, with fixed-prefix stability.  Transcendence does not pass through limits by itself.\\[6pt]
\hyperref[mi:gap-theorem]{M4} & Divisible prime gaps and a transcendental gap exponent & Necessary consequences of hypothetical algebraicity of $M$.  Their consistency has not been ruled out.\\[6pt]
\hyperref[mi:unit-polynomial]{M5} & An exact three-trace elimination polynomial & Proved on the unit-norm branch.  It cannot be applied to every algebraic candidate without proving that extra restriction.\\[6pt]
\hyperref[mi:depth]{M6} & Three-adic depth and prime splitting restrictions & Exact source hypotheses are retained.  The conclusions do not cover arbitrary cubic Pisot numbers.\\[6pt]
\hyperref[mi:seed]{M7} & A prime seed with no later divisibility return & Explicit polynomial and exact finite-state certificate of a period-$85\,712$ orbit.  This is not an all-prime infinite sequence.\\[6pt]
\hyperref[mi:preserve]{M8} & Quantitative conflict between prefix preservation and a shift-escape certificate & A limitation of the displayed sufficient estimates, not a universal impossibility theorem for perturbation methods.\\[6pt]
\hyperref[mi:terminal]{M9} & Collapse of a skipped-level corridor to the last adjacent corridor & All intermediate prime constraints are retained.  A larger raw interval is not automatically an admissible lower branch.\\[6pt]
\hyperref[mi:rh]{M10} & DH/RH conditional transcendence and the scope of a quasi-Riemann input & External conditional theorem and a scope comparison.  No Density Hypothesis or RH is proved in this manuscript.\\
\bottomrule
\end{longtable}
\normalsize
\clearpage

\section{Winding embeddings and arithmetic fields}
\label{gf:section}

This section gives the structural results behind entries G1--G5 in the
concordance.  Its external analytic input is the canonical phase
decomposition of $E$-functions and the value morphism of
Fischler--Rivoal \cite[Theorem~2.3 and Proposition~3.1]{FR2026}, together
with Beukers' division theorem \cite{Beukers2006}.  The continuation
argument and the finite-phase method were developed in the working
manuscripts \cite{VMRC030,VMRC033,VMRC049}; the proofs needed here are
included below.  These are statements about specified subfields and
subrings of $E$-values.  They neither identify the arithmetic nature of
$\gamma$ nor assert historical priority for their formulations.

\subsection{The ambient ring and the analytic input}
\label{gf:ambient}

Throughout this section $\A=\overline{\mathbb Q}\subset\mathbb C$.
Let $\Er$ be the ring of values at algebraic points of $E$-functions, and
put $\K=\operatorname{Frac}(\Er)$.  One may use the usual exponential
arithmetic bounds in the definition of an $E$-function: if
$f(z)=\sum_{n\ge0}a_nz^n/n!$, its coefficients are algebraic, their
conjugates and common denominators have exponential bounds, and $f$
satisfies a linear differential equation over $\A(z)$.  All the explicit
functions below belong to this class.  The construction also applies to
the broader Siegel convention covered in \cite{FR2026}.

Every element of $\Er$ can be written as $f(1)$: algebraic changes of
scale and finite sums and products preserve the class of $E$-functions.
In particular, $\A\subset\Er$ and $e^a,e^{-a}\in\Er$ for $a\in\A$.
Define
\begin{equation}\label{gf:Ein-def}
 \Ein(z)=\sum_{n\ge1}\frac{(-1)^{n-1}z^n}{n\,n!},
 \qquad \eta_a=\Ein(a),\qquad \eta=\eta_1.
\end{equation}
For algebraic $a$, the function $\Ein(az)$ is an $E$-function.  Its
series has the required arithmetic bounds, and it satisfies
$z f'''+(2+az)f''+af'=0$.
For real $a>0$, put
\[
 \Eone(a)=\int_a^\infty\frac{e^{-t}}{t}\,dt,
 \qquad \delta=e\Eone(1)=\int_0^\infty\frac{e^{-t}}{1+t}\,dt.
\]
The classical connection identity gives
\begin{equation}\label{gf:connection}
 \eta_a=\gamma+\log a+\Eone(a),
 \qquad \delta=e(\eta-\gamma).
\end{equation}
No membership of $\gamma$ or $\delta$ in $\Er$ is assumed here.

We shall use the ring of finite exponential sums
\begin{equation}\label{gf:R}
 \mathcal R=\left\{
       \sum_{\rho\in\A}c_\rho e^{\rho X}:
       c_\rho\in\mathbb C,\quad\text{finite support}
                         \right\}.
\end{equation}
This is a domain of entire functions.  Distinct exponential functions
are linearly independent over $\mathbb C$: differentiating a finite
linear relation at $X=0$ gives a Vandermonde system.

Fix $-\pi/2<\theta<\pi/2$.  The cited phase decomposition associates to
an $E$-function $f$ finitely many germs $f_\rho$ at $z=1$ with
\begin{equation}\label{gf:phase-decomp}
 f(z)=\sum_\rho f_\rho(z)e^{\rho z}.
\end{equation}
They are the canonical sums of the corresponding asymptotic components.
At an exceptional direction one uses $\theta+\epsilon$ with
$\epsilon>0$ sufficiently small.  For any finite collection of
functions, a common such $\epsilon$ is available, and the germs are
independent of sufficiently small positive changes of $\epsilon$.
The following properties of the construction are the external input:
the decomposition is unique; sums and products have their corresponding
componentwise and convolution decompositions; and each
$e^{\rho z}f_\rho(z)$ solves an $E$-operator with no possible finite
singularity except $0$.  In particular, the germs $f_\rho$ continue along
every path in $\mathbb C^\times$.  The last property, stated just before
Theorem~2.3 of \cite{FR2026}, is essential for the argument below.

\begin{proposition}[Continued phase embeddings]\label{gf:windings}
Let $\mathsf M^j$ mean analytic continuation along a based loop at $1$
with counterclockwise winding number $j\in\mathbb Z$ about $0$.
For every integer $j$, the prescription
\begin{equation}\label{gf:Phi-def}
 \Phi_\theta^{[j]}(f(1))(X)
       =\sum_\rho(\mathsf M^j f_\rho)(1)e^{\rho X}
\end{equation}
defines an injective $\A$-algebra homomorphism $\Er\longrightarrow
\mathcal R$.  It satisfies
\begin{equation}\label{gf:Phi-normalization}
 \Phi_\theta^{[j]}(\xi)(1)=\xi,
 \qquad \Phi_\theta^{[j]}(e^a)=e^{aX}\quad(a\in\A).
\end{equation}
It extends uniquely to an embedding
$\K\longrightarrow\operatorname{Frac}(\mathcal R)$.
\end{proposition}
\begin{proof}
Suppose $f(1)=g(1)$.  Beukers' division theorem supplies an $E$-function
$h$ such that $f-g=(z-1)h$.  Choose a common sufficiently small lateral
direction for $f,g,h$.  Uniqueness of their canonical decompositions
implies
\[
 f_\rho-g_\rho=(z-1)h_\rho
\]
for each phase, with missing components interpreted as zero.  Continue
this identity along the same based loop.  The continued $h_\rho$ is
holomorphic at the endpoint $1$, since $1$ is not a singularity of its
differential equation.  Evaluation there yields
$(\mathsf M^j f_\rho)(1)=(\mathsf M^j g_\rho)(1)$.  Thus
\eqref{gf:Phi-def} is independent of the representing $E$-function.

The component identities for sums give additivity.  For a product
$h=fg$, canonical uniqueness gives
\[
 h_\rho=\sum_{\alpha+\beta=\rho}f_\alpha g_\beta.
\]
Continuation respects products, so this identity proves multiplicativity
after evaluation at $1$.  Constants and pure exponentials have constant
components and give the claimed normalizations.  Finally, the total
function $f$ is entire and hence single-valued.  Continuing
\eqref{gf:phase-decomp} gives
\[
 \sum_\rho(\mathsf M^j f_\rho)(1)e^\rho
       =(\mathsf M^j f)(1)=f(1).
\]
Evaluation at $X=1$ is therefore a left inverse, proving injectivity.
The extension to fractions is now immediate.  Although evaluation is
not defined on every element of $\operatorname{Frac}(\mathcal R)$ at
$1$, it is defined on this embedded copy of $\K$: a denominator
$\Phi_\theta^{[j]}(b)$, with $0\ne b\in\Er$, takes the nonzero value
$b$ at $1$.
\end{proof}

\begin{remark}[What is continued]\label{gf:what-winds}
The entire function $f$ and its value $f(1)$ do not acquire monodromy.
Only its separate canonical components are continued.  Nor is
$\Phi_\theta^{[j]}$ defined by replacing $\theta$ with $\theta+2\pi j$:
the initial summation direction stays fixed.  The continued components
need not constitute a canonical asymptotic decomposition in a specified
new sector.  These distinctions prevent a change of analytic branch
from being mistaken for a change of the number $\gamma$.
\end{remark}

\subsection{The affine winding formula for exponential integrals}
\label{gf:Ein-winding}

Let $a>0$ be algebraic and choose the positive real normalization.
Writing $\tau=2\pi i$, the two components of $f(z)=\Ein(az)$ are
\[
 f_0(z)=\gamma+\log a+\log z,
 \qquad f_{-a}(z)=e^{az}\Eone(az).
\]
They have phases $0$ and $-a$.  Their canonical identification follows
either from the connection identity or from the inverse-power expansion
of $e^{az}\Eone(az)$; its Borel transform is $1/(a+t)$, which has no
singularity on the positive ray.  Since $f$ is entire, continuation
around $0$ adds $j\tau$ to $f_0$ and subtracts $j\tau e^{az}$ from
$f_{-a}$.  Proposition~\ref{gf:windings} therefore gives
\begin{equation}\label{gf:Ein-winding-formula}
 \Phi_0^{[j]}(\eta_a)
   =C_{a,j}+B_{a,j}e^{-aX},
 \quad
 \begin{cases}
 C_{a,j}=\gamma+\log a+j\tau,\\
 B_{a,j}=e^a\Eone(a)-j\tau e^a.
 \end{cases}
\end{equation}
At $j=1$, $B_{a,1}\ne0$ because $e^a\Eone(a)>0$ is real.  Also
\begin{equation}\label{gf:trans-centers}
 C_{a,1}\notin\A.
\end{equation}
Indeed its imaginary part is $2\pi$.  The imaginary part of an algebraic
complex number is algebraic, whereas $\pi$ is transcendental.  This is
where positivity of $a$ supplies the reality needed for the joint
theorem; no assertion about the arithmetic of $\gamma+\log a$ is used.

\subsection{Finite phases and ramification}
\label{gf:descent}

\begin{lemma}[Finite-phase descent]\label{gf:phase-descent}
Let $a_1,\ldots,a_r\in\A$ be linearly independent over $\mathbb Q$.
If $W\in\operatorname{Frac}(\mathcal R)$ is algebraic over
$\mathbb C(e^{a_1X},\ldots,e^{a_rX})$, then for some integer $q\ge1$,
\begin{equation}\label{gf:finite-phase-conclusion}
 W\in\mathbb C(e^{a_1X/q},\ldots,e^{a_rX/q}).
\end{equation}
If in addition $W\in\mathcal R$, it lies in the Laurent polynomial
ring in these variables after a possible increase of $q$.
\end{lemma}
\begin{proof}
Choose a rational basis extending $a_1,\ldots,a_r$ for the span of all
phases in a numerator and denominator representing $W$.  Clearing
their rational coordinates gives variables
$X_i=e^{a_iX/q}$ and additional variables $Y_j=e^{b_jX/q}$, with
\[
 W\in\mathbb C(X_1,\ldots,X_r,Y_1,\ldots,Y_s).
\]
These variables are algebraically independent over $\mathbb C$:
a polynomial relation would be a linear relation among exponentials
with distinct exponents.  Put $F_0=\mathbb C(X_1,\ldots,X_r)$.
The element $W$ is algebraic over $F_0$, because it is algebraic over
the smaller field $\mathbb C(X_1^q,\ldots,X_r^q)$.

A purely transcendental extension $F_0(Y_1,\ldots,Y_s)$ has no new
elements algebraic over $F_0$.  To check this directly, differentiate
the separable minimal polynomial of such an element with respect to
each $Y_j$.  Its derivative at the element is nonzero, so each partial
derivative of the element vanishes.  A rational function in
characteristic zero with zero derivative in a variable is independent
of that variable.  This proves \eqref{gf:finite-phase-conclusion}.
In the ring case, $W$ is initially a Laurent polynomial in the $X_i,Y_j$;
uniqueness of its Laurent coefficients removes every nonconstant
$Y$-monomial.
\end{proof}

\begin{lemma}[Preservation of degree under constants]\label{gf:constant-descent}
Let $U_1,\ldots,U_r$ be independent indeterminates.  An irreducible
polynomial in $\A(U_1,\ldots,U_r)[T]$ stays irreducible over
$\mathbb C(U_1,\ldots,U_r)$.  Consequently a finite extension of
$\A(U_1,\ldots,U_r)$ retains its degree when constants are extended
to $\mathbb C$.
\end{lemma}
\begin{proof}
By clearing denominators and taking a primitive part, Gauss's lemma
reduces the assertion to an irreducible polynomial over the
algebraically closed field $\A$.  If it factored after extending
constants, factors of bounded degree with suitable normalized
coefficients would satisfy finitely many polynomial equations over
$\A$.  Nonzero normalizing coefficients can be enforced by adjoining
their inverses.  A solution over $\mathbb C$ would then give a solution
over $\A$ by the weak Nullstellensatz, contradicting irreducibility.
Applying this observation to the polynomial of a primitive element
proves the degree statement.  The same argument applies after any
identification of the independent $U_i$ with algebraically independent
functions over $\mathbb C$.
\end{proof}

\begin{lemma}[Coordinate inertia in a translated Kummer extension]
\label{gf:Kummer-inertia}
Let $C_i,B_i\in\mathbb C$, $B_i\ne0$, and let $Z_1,\ldots,Z_r$
be independent variables with $U_i=C_i+B_iZ_i^q$.  The extension
\[
 \mathbb C(Z_1,\ldots,Z_r)/\mathbb C(U_1,\ldots,U_r)
\]
is Galois with group $G=(\mathbb Z/q\mathbb Z)^r$.
At the generic point of $U_i=C_i$ its inertia group is precisely the
$i$th coordinate subgroup $G_i$.
If an intermediate extension is unramified at every such hyperplane,
it is the base field.
\end{lemma}
\begin{proof}
Each coordinate transformation multiplies $Z_i$ by an independent
$q$th root of unity, and the degree is $q^r$.  Along $Z_i=0$ the
function $U_i-C_i=B_iZ_i^q$ has order $q$.  The other variables are
units, and the transformations of $Z_j$ for $j\ne i$ act nontrivially
on the residue field.  The transformations of $Z_i$ act trivially on
that residue field, giving inertia exactly $G_i$.

Let $H\subseteq G$ correspond to an intermediate field.  Since $G$
is abelian, its inertia at this divisor is the image of $G_i$ in
$G/H$.  If all these images vanish, all $G_i$ lie in $H$.  They generate
$G$, so $H=G$ and the intermediate field is the base field.
\end{proof}

The coordinate divisors here are valuations of abstract rational
function fields.  Using them does not require an exponential function
to take the value zero at a finite value of $X$.

\begin{theorem}[Relative algebraic closure for positive Ein families;
G1]\label{gf:jointclosure}
Let $a_1,\ldots,a_r>0$ be algebraic numbers linearly independent over
$\mathbb Q$, and put
\[
 F=\A(\eta_{a_1},\ldots,\eta_{a_r}).
\]
The generators $\eta_{a_1},\ldots,\eta_{a_r}$ are algebraically
independent, and $F$ is relatively algebraically closed in $\K$:
\begin{equation}\label{gf:jointclosure-equation}
 \{v\in\K:v\text{ is algebraic over }F\}=F.
\end{equation}
\end{theorem}
\begin{proof}
Use $\Phi=\Phi_0^{[1]}$ and write
$C_i=C_{a_i,1}$, $B_i=B_{a_i,1}$.  Its images
\[
 U_i=\Phi(\eta_{a_i})=C_i+B_ie^{-a_iX}
\]
are algebraically independent over $\mathbb C$ by the rational
independence of the $a_i$ and the exponential Vandermonde argument.
A polynomial relation over $\A$ among the actual values would give
one among these functions.  Thus the actual values are algebraically
independent over $\A$ as well.

Let $v\in\K$ be algebraic over $F$, of degree $d$.  Its image
$W=\Phi(v)$ is algebraic over $\mathbb C(U_1,\ldots,U_r)$.
Finite-phase descent supplies an integer $q$ for which
\[
 W\in\mathbb C(Z_1,\ldots,Z_r),\qquad
 Z_i=e^{-a_iX/q},\qquad U_i=C_i+B_iZ_i^q.
\]
The degree of $\mathbb C(U,W)/\mathbb C(U)$ is $d$ by
Lemma~\ref{gf:constant-descent}.

Take a primitive irreducible polynomial
$P(U,T)\in\A[U_1,\ldots,U_r,T]$ for $v$ after denominators are cleared.
Its leading coefficient in $T$ times its $T$-discriminant is a nonzero
polynomial $D(U)\in\A[U_1,\ldots,U_r]$.  Outside $D=0$, the
corresponding algebraic extension is unramified.  For example, after
inverting the leading coefficient, the polynomial is monic, and
inverting the discriminant makes the algebra generated by a root
finite etale.  Thus every ramified divisor in the extension must lie
in $D=0$.

Suppose that the coordinate inertia at $U_i=C_i$ survived in
$\mathbb C(U,W)$.  The whole hyperplane would then be contained in
$D=0$.  Expand $D$ as a polynomial in the variables other than $U_i$.
At least one coefficient is a nonzero polynomial in $\A[U_i]$,
and its value at $C_i$ would be zero.  This would make $C_i$ algebraic,
contrary to \eqref{gf:trans-centers}.  No coordinate inertia can
survive.  Lemma~\ref{gf:Kummer-inertia} implies
$\mathbb C(U,W)=\mathbb C(U)$, so $d=1$.  Therefore $v\in F$.
The reverse inclusion in \eqref{gf:jointclosure-equation} is immediate.
\end{proof}

\begin{corollary}[No radical of a shifted Ein value]\label{gf:Ein-radicals}
For positive algebraic $a$, algebraic $c$, and an integer $n\ge2$, no
solution of
\[
 w^n=\eta_a-c
\]
belongs to $\K$.  In particular $\eta^{1/n}\notin\K$.
\end{corollary}
\begin{proof}
Such a $w$ would be algebraic over $\A(\eta_a)$ and hence belong to
that field by Theorem~\ref{gf:jointclosure}.  In the rational function
field $\A(U)$, the valuation of $U-c$ at $U=c$ equals $1$, whereas
the valuation of an $n$th power is divisible by $n$.  This is impossible.
\end{proof}

\subsection{Rational intersections and the exponential--Ein pair}
\label{gf:intersections}

\begin{lemma}[Two affine images exclude rational denominators]
\label{gf:two-centers-ring}
Suppose $\xi\in\Er$ has, under two injective $\A$-algebra
homomorphisms $\Psi_1,\Psi_2:\Er\to\mathcal R$, images
\[
 \Psi_i(\xi)=C_i+B_i e^{a_iX},\qquad
 B_i a_i\ne0,\quad C_1\ne C_2.
\]
Then $\xi$ is transcendental and
\begin{equation}\label{gf:ring-intersection-general}
 \A(\xi)\cap\Er=\A[\xi],\qquad
 \A[\xi]\cap\Er^\times=\A^\times.
\end{equation}
\end{lemma}
\begin{proof}
A nonconstant image rules out $\xi\in\A$, since both maps fix $\A$.
Let $P,Q\in\A[T]$ be coprime and suppose
$v=P(\xi)/Q(\xi)\in\Er$.  Applying $\Psi_i$ to
$Q(\xi)v=P(\xi)$ gives an identity of entire functions.  If $\alpha$
were a root of $Q$ distinct from $C_i$, some complex $X_0$ would
satisfy $e^{a_iX_0}=(\alpha-C_i)/B_i$.  Evaluation there would force
$P(\alpha)=0$, contradicting coprimality.  Hence every root of $Q$
would have to equal both $C_1$ and $C_2$.  Thus $Q$ is constant.
This proves the first assertion; its reverse inclusion is immediate.
If a polynomial in $\xi$ is a unit of $\Er$, its reciprocal is another
polynomial in $\xi$ by the first assertion.  The only units of the
polynomial ring $\A[\xi]$ are the nonzero constants, proving the second.
\end{proof}

\begin{theorem}[The rational Ein intersection; G4]\label{gf:einintersection}
For every nonzero algebraic $a$,
\begin{equation}\label{gf:einintersection-equation}
 \A(\eta_a)\cap\Er=\A[\eta_a],\qquad
 \A[\eta_a]\cap\Er^\times=\A^\times.
\end{equation}
In particular, $\eta_a$ is transcendental and
$1/\eta_a\notin\Er$.
\end{theorem}
\begin{proof}
For positive $a$ this follows directly from
\eqref{gf:Ein-winding-formula} and Lemma~\ref{gf:two-centers-ring},
using any two windings with nonzero exponential coefficient.
For completeness the same argument works at complex algebraic $a$,
with the following branch bookkeeping.  Choose one admissible
canonical direction for $\Ein(az)$.  Its two phases are $0,-a$, and
its components have the form
\[
 f_0(z)=C+\log z,\qquad
 f_{-a}(z)=e^{az}\bigl(\Ein(az)-C-\log z\bigr),
\]
where $C$ is fixed by the chosen canonical sum and the compatible
logarithm.  This is the connection identity analytically continued
from a nonsingular summation ray.  Equivalently, the inverse-power
component has Borel transform $1/(a+t)$; if a ray meets its pole,
one fixes either adjacent nonsingular ray and the compatible branch.
Only the existence of this constant $C$, not its arithmetic nature,
is needed here.

Write $B=e^a(\eta_a-C)$.  Since $\Ein(az)$ and $e^{az}$ are
single-valued, continuing the components gives
\[
 \Phi_\theta^{[j]}(\eta_a)
    =C+j\tau+(B-j\tau e^a)e^{-aX}.
\]
The coefficient $B-j\tau e^a$ vanishes for at most one integer $j$.
Choose two distinct windings avoiding it.  The omitted values
$C+j\tau$ differ, so Lemma~\ref{gf:two-centers-ring} applies.
\end{proof}

The assertion $1/\eta_a\notin\Er$ must not be upgraded to
$1/\eta_a\notin\K$.  In fact $1/\eta_a\in\K$ by the definition of
a fraction field.  The radical exclusion in
Corollary~\ref{gf:Ein-radicals}, in contrast, is genuinely a statement
about the larger field $\K$.

\begin{theorem}[The exponential--Ein pair; G5]\label{gf:epair}
The numbers $e$ and $\eta$ are algebraically independent over $\A$.
If $\gamma\in\A$, then $\delta$ is transcendental over $\A(\eta)$.
\end{theorem}
\begin{proof}
Put $C=\gamma+\tau$ and $B=\delta-\tau e$.  Under $\Phi_0^{[1]}$,
\[
 e\longmapsto e^X,
 \qquad \eta\longmapsto C+Be^{-X}.
\]
Here $C\notin\A$ by \eqref{gf:trans-centers}.  Suppose that
$\sum_{r=0}^d e^r P_r(\eta)=0$, with $P_r\in\A[T]$ and
$P_d\ne0$.  Apply the embedding and divide by $e^{dX}$.  On letting
real $X\to+\infty$, all terms with $r<d$ tend to zero, while the
remaining term tends to $P_d(C)$.  This is nonzero because $C$ is
transcendental.  The contradiction proves algebraic independence.

If $\gamma=q\in\A$, then $\delta=e(\eta-q)$ and $\eta-q\ne0$.
Algebraicity of $\delta$ over $\A(\eta)$ would imply algebraicity of
$e=\delta/(\eta-q)$ over that field, contradicting the first assertion.
\end{proof}

The independence of this pair is an application of the established
$E$-function framework, not an assertion that the general
Siegel--Shidlovskii phenomenon is new here.  Its role is to prevent the
conditional identity $\delta=e(\eta-q)$ from being mistaken for a
contradiction: if $\gamma=q$ were algebraic, the required transcendence
of $\delta$ over $\A(\eta)$ would be mathematically consistent with
the preceding statements.

\subsection{A radical criterion for Euler's constant}
\label{gf:radical-criterion}

\begin{theorem}[The Euler--Gompertz radical obstruction; G2]
\label{gf:delta-roots}
For every integer $n\ge2$,
\begin{equation}\label{gf:delta-roots-equation}
 \gamma\in\A\quad\Longrightarrow\quad\delta^{1/n}\notin\K.
\end{equation}
Consequently, an independent proof that $\delta^{1/n}\in\K$ for
some $n\ge2$ would prove that $\gamma$ is transcendental.
\end{theorem}
\begin{proof}
Assume $\gamma=q\in\A$ and that a root $v$ of $v^n=\delta$ belongs
to $\K$.  The number $e^{1/n}$ is an $E$-value, so
$w=v/e^{1/n}\in\K$.  Equation~\eqref{gf:connection} gives
\[
 w^n=\eta-q,
\]
contrary to Corollary~\ref{gf:Ein-radicals}.  Different choices of the
root differ by an algebraic root of unity, so the conclusion is
independent of that choice.  The final assertion is the contrapositive.
\end{proof}

\begin{remark}[The missing arithmetic membership]\label{gf:missing-membership}
No inclusion $\delta^{1/n}\in\K$ is proved in this paper.  In
particular, the theorem is not an irrationality proof for $\gamma$.
The same argument shows, for every $q\in\A$, that
\[
 \bigl(e(\eta-q)\bigr)^{1/n}\notin\K\qquad(n\ge2).
\]
Thus the exclusion alone does not distinguish the actual Euler
constant from a hypothetical algebraic parameter in the identity
$e(\eta-q)$.  To use the criterion, the additional membership must
come from an independent property of the actual constant $\delta$.
\end{remark}

\subsection{A compositum counterexample}
\label{gf:compositum-section}

We first record a one-variable variant of the finite-cover argument.
Unlike Theorem~\ref{gf:jointclosure}, it does not require either
omitted value to be transcendental.

\begin{lemma}[Two affine images exclude finite algebraic extensions]
\label{gf:two-centers-field}
Let $\xi\in\Er$.  Suppose two embeddings of the form in
Proposition~\ref{gf:windings} give
\[
 \Psi_i(\xi)=C_i+B_i e^{a_iX},
 \qquad a_i\in\A^\times,\quad B_i\ne0,\quad C_1\ne C_2.
\]
Then $\A(\xi)$ is relatively algebraically closed in $\K$.
\end{lemma}
\begin{proof}
The nonconstant images imply that $\xi$ is transcendental.  Let
$L=\A(\xi,v)$ for an element $v\in\K$ algebraic of degree $d$ over
$\A(\xi)$.  Identify the transcendental generator $\xi$ with an
indeterminate $U$.  After extending constants to $\mathbb C$, this
gives a function field $L_{\mathbb C}$ of degree $d$ over
$\mathbb C(U)$ by Lemma~\ref{gf:constant-descent}.

For either embedding, finite-phase descent places the image of $v$
in $\mathbb C(Z_i)$ with $Z_i=e^{a_iX/q_i}$ and
$U=C_i+B_iZ_i^{q_i}$ for some $q_i\ge1$.  The comparison is made over
the same abstract coordinate $U$; it does not identify the two
functions $\Psi_i(v)$ inside their ambient exponential fields.

If $d>1$, the first realization must be ramified at $U=C_1$.
Indeed, in the one-coordinate case of
Lemma~\ref{gf:Kummer-inertia}, the inertia group at that point is
the entire cyclic group, and its image in every nontrivial intermediate
extension is nontrivial.  On the other hand, the second realization
is unramified at every finite point other than $C_2$.  One may see
this directly by solving
$Z_2^{q_2}=(U-C_2)/B_2$ locally at a nonzero right-hand side, where its
roots are simple.  Since $C_1\ne C_2$, the same abstract extension
would be both ramified and unramified at $C_1$, a contradiction.
Therefore $d=1$ and $v\in\A(\xi)$.
\end{proof}

\begin{theorem}[Relative closure is not preserved by composita; G3]
\label{gf:compositum}
For every $q\in\A$ put
\[
 d_q=e(\eta-q),\qquad F_q=\A(d_q).
\]
Each $F_q$ is relatively algebraically closed in $\K$.  If $q\ne r$,
however, then
\begin{equation}\label{gf:compositum-equation}
 F_qF_r=\A(e,\eta),
\end{equation}
and this compositum is not relatively algebraically closed in $\K$.
More precisely, $e^{1/n}\in\K\setminus\A(e,\eta)$ for every $n\ge2$.
\end{theorem}
\begin{proof}
All $d_q$ belong to $\Er$.  Formula~\eqref{gf:Ein-winding-formula}
at $a=1$ gives, without assuming $\gamma\in\Er$,
\begin{equation}\label{gf:dq-images}
 \Phi_0^{[j]}(d_q)
   =(\gamma-q+j\tau)e^X+\delta-j\tau e.
\end{equation}
In this formula $\gamma$ and $\delta$ are complex coefficients of
the image, not additional presumed elements of the domain of $\Phi$.
The coefficient of $e^X$ vanishes for at most one integer $j$.
Choose two distinct other windings.  Their constant terms
$\delta-j\tau e$ differ.  Lemma~\ref{gf:two-centers-field} proves
relative algebraic closure of $F_q$.

For distinct $q,r$, the identities
\[
 e=\frac{d_q-d_r}{r-q},\qquad
 \eta=q+\frac{d_q}{e}
\]
prove \eqref{gf:compositum-equation}.  The reverse containment follows
from the definitions of $d_q,d_r$.
By Theorem~\ref{gf:epair}, $\A(e,\eta)$ is a rational function field
in the two independent generators $e,\eta$.  The polynomial
$T^n-e$ is irreducible over this field: view it as $T^n-U$ over
$\A(V)(U)$ and apply Eisenstein's criterion at the prime $U$ in
$\A(V)[U]$.  Thus $e^{1/n}\notin\A(e,\eta)$, although
$e^{1/n}\in\Er\subset\K$.  This proves the failure of relative
algebraic closure of the compositum.
\end{proof}

The same affine images, by Lemma~\ref{gf:two-centers-ring}, also give
$\A(d_q)\cap\Er=\A[d_q]$ and
$\A[d_q]\cap\Er^\times=\A^\times$.
If $\gamma=q$ were algebraic, then $\delta=d_q$ would be a
transcendental nonunit of $\Er$, consistently with these statements.
The compositum example is a concrete reason that separate closure
theorems cannot be combined without a new proof of closure for their
jointly generated field.

\begin{remark}[Differences of phase images are not new $E$-values]
\label{gf:no-pullback}
For $j\ne0$,
\[
 \Phi_0^{[j]}(\eta)-\Phi_0^{[0]}(\eta)
     =j\tau(1-e e^{-X})\ne0.
\]
This difference vanishes at $X=1$.  It cannot lie in the image of any
continued phase embedding of an $E$-value: evaluation would force
that value to be $0$, whose image is identically zero.  Subtracting
two different embeddings therefore does not provide an inverse-image
construction for $\pi$, $\gamma$, or $\delta$.  All membership
assertions used in this section were established in the original
ring or field before applying an embedding.
\end{remark}

\section{Shifted prime clocks and fixed polynomial relations}
\label{mi:shifts}

For a positive real sequence \(C=(C_k)_{k\ge1}\), define
\[
 \mathcal W(C)=
 \{A>1:\floor{A^{C_k}}\text{ is prime for every }k\ge1\}.
\]
When this set is nonempty, its least element is denoted by \(\xi(C)\).
The starting index is part of this definition. In particular,
\[
 M=\xi((3^k)_{k\ge1}),\qquad
 A_\pm=\xi((3^k\pm1)_{k\ge1}),\qquad
 B_\pm=\xi((3^k\pm k)_{k\ge1}).
\]
Existence for the clocks used below follows from Saito's existence
theorem, since their successive exponent ratios are at least
\(2\) \cite[Theorem~1.3]{SaitoVariants}.
The irrationality of the original \(M\) is an external result
\cite{Saito2025}; its unconditional transcendence is not proved here.

\paragraph{Status of the named shifted examples.}
The transcendence statements for \(A_\pm\) and \(B_\pm\) were
inputs to the research discussion. Their source proofs are in the
internal manuscripts \cite{VMRC064,VMRC008}, with additional
extensions and audits in \cite{VMRC075}. We do not reproduce or
independently certify those complete proofs in this paper. The
theorems below are therefore formulated for an arbitrary
\emph{transcendental} real base \(A>1\); their proofs are independent
of the status of those four particular examples. Primality is
unnecessary for the polynomial-avoidance assertions themselves.

\subsection{Two coordinates for a constant shift}
\label{mi:two}

\begin{theorem}[Fixed polynomial avoidance for constant shifts]
\label{mi:affine}
Let \(A>1\) be transcendental, let \(\sigma\) be a nonzero integer,
and assume \(3^k+\sigma>0\) at the indices under consideration.
Put \(p_k=\floor{A^{3^k+\sigma}}\). For every fixed \(h\ge1\),
\[
 \frac{p_{k+h}}{p_k^{3^h}}
 \longrightarrow A^{\sigma(1-3^h)}.
\]
For every nonzero \(P\in\A[X,Y]\),
\[
 P(p_k,p_{k+h})\ne0
 \qquad\text{for all sufficiently large }k.
\]
The threshold may depend on \(A,\sigma,h,P\).
\end{theorem}

\begin{proof}
Write \(d=3^h\) and
\(L=A^{\sigma(1-d)}\). The elementary estimate
\[
 p_k=A^{3^k+\sigma}\bigl(1+O(A^{-3^k-\sigma})\bigr)
\]
proves the ratio limit. Since \(\sigma(1-d)\) is a nonzero integer,
\(L\) is transcendental. Write
\(P(X,Y)=\sum c_{ij}X^iY^j\) and let
\(D=\max\{i+dj:c_{ij}\ne0\}\). Then
\[
 \frac{P(p_k,p_{k+h})}{p_k^D}
 \longrightarrow \sum_{i+dj=D}c_{ij}L^j.
\]
For a given \(j\), the equation \(i+dj=D\) specifies at most one
\(i\). The polynomial on the right is consequently nonzero, and
cannot vanish at the transcendental number \(L\).
\end{proof}

For the two constant shifts in the comparison register, this gives
\[
 \frac{p_{+,k+1}}{p_{+,k}^3}\longrightarrow A_+^{-2},
 \qquad
 \frac{p_{-,k+1}}{p_{-,k}^3}\longrightarrow A_-^{2}.
\]
The result excludes a \emph{fixed} algebraic-coefficient recurrence
on infinitely many pairs. It does not say that an individual pair
of integers is algebraically independent.

\subsection{Three coordinates for a linear shift}
\label{mi:three}

\begin{theorem}[Fixed polynomial avoidance for linear shifts]
\label{mi:linear}
Let \(A>1\) be transcendental and \(\varepsilon\in\{1,-1\}\).
Put \(q_k=\floor{A^{3^k+\varepsilon k}}\). Then
\[
 \frac{q_{k+2}q_k^3}{q_{k+1}^4}
 \longrightarrow A^{-2\varepsilon}.
\]
For every nonzero \(P\in\A[X,Y,Z]\),
\[
 P(q_k,q_{k+1},q_{k+2})\ne0
 \qquad\text{for all sufficiently large }k.
\]
\end{theorem}

\begin{proof}
The exponent in the ratio is exactly
\[
 (3^{k+2}+\varepsilon(k+2))
 +3(3^k+\varepsilon k)
 -4(3^{k+1}+\varepsilon(k+1))=-2\varepsilon,
\]
which proves the first claim.

For a monomial \(X^iY^jZ^\ell\), define
\[
 W=i+3j+9\ell,\qquad V=i+j+\ell,\qquad U=j+2\ell.
\]
Its unfloored value is \(A^{W3^k+\varepsilon Vk+\varepsilon U}\).
The linear map \((i,j,\ell)\mapsto(W,V,U)\) is injective because
\[
 \det\begin{pmatrix}1&3&9\\1&1&1\\0&1&2\end{pmatrix}=4.
\]
Among the finitely many monomials of \(P\), first maximize \(W\),
then maximize \(\varepsilon V\). In this leading group, distinct
monomials have distinct \(U\), so their remaining coefficient is
a nonzero Laurent polynomial in \(A\), with algebraic coefficients.
It is nonzero by transcendence of \(A\).

Terms with smaller \(W\) are exponentially smaller on the
\(3^k\) scale; terms with the same \(W\) but smaller
\(\varepsilon V\) are exponentially smaller on the \(k\) scale.
Replacing unfloored powers by floors has relative error
\(O(A^{-3^k+O(k)})\) in every fixed monomial. After normalization
by the leading \(A^{W3^k+\varepsilon Vk}\), these errors also
tend to zero. The normalized polynomial therefore tends to the
nonzero coefficient just identified.
\end{proof}

The extra coordinate is explained by the three scales
\(3^k,k,1\). These theorems apply separately to each chosen base.
They do not assert joint algebraic independence of
\(A_+,A_-,B_+,B_-\).

\subsection{Continuity at the original clock}
\label{mi:continuity}

The following continuity theorem is an analytic connection with
the actual minimum; it is not an arithmetic limit theorem.
The argument, recorded in \cite{VMRC075}, is included because
finite initial constraints cannot be discarded in this problem.

\begin{lemma}[Uniform continuation from a large prime]
\label{mi:continuation}
There is \(P_0\) such that every prime \(p\ge P_0\) at any level
of a clock whose successive ratios lie in \([5/2,4]\) has an
infinite prime continuation with strictly nested closed
inverse-power cells.
\end{lemma}

\begin{proof}
The Baker--Harman--Pintz bound supplies a constant \(K\) such that
\(\operatorname{nextprime}(x)-x\le Kx^{21/40}\) for all sufficiently
large \(x\) \cite{BHP2001}. Uniformly for \(5/2\le c\le4\),
\[
 (p+1)^c-p^c\ge cp^{c-1},\qquad
 \frac{cp^{c-1}}{p^{21c/40}}
 =cp^{19c/40-1}\ge\frac52p^{3/16}.
\]
Hence, for all sufficiently large \(p\), a prime \(q\) can be
chosen with \(p^c<q\) and \(q+1<(p+1)^c\).
Iteration gives strictly nested closed root cells, whose lengths
tend to zero. Their common point lies below each upper endpoint,
because the next cell already does, and thus has the prescribed
prime floors.
\end{proof}

\begin{theorem}[Continuity and normalized tails]
\label{mi:tail-limit}
For \(-1\le t\le1\), put
\[
 F(t)=\xi((3^k-t)_{k\ge1}).
\]
Then \(F(t)\to M\) as \(t\to0\). Consequently,
\[
 \xi((3^{k+h}-1)_{k\ge1})^{3^h}\longrightarrow M,\qquad
 \xi((3^{k+h}+1)_{k\ge1})^{3^h}\longrightarrow M.
\]
At each fixed finite set of indices, the prime outputs of these
normalized minima on their own perturbed clocks eventually
coincide with the outputs of \(M\).
\end{theorem}

\begin{proof}
The ratios \((3^{k+1}-t)/(3^k-t)\) lie in \([5/2,4]\).
Lemma~\ref{mi:continuation}, starting at one fixed sufficiently
large first prime, gives a uniform upper bound on \(F(t)\).
The first floor gives \(F(t)\ge2^{1/4}>1\).

Suppose \(t_n\to0\) and \(F(t_n)\to A\). If \(A^{3^j}=m\)
were an integer, \(m\ge2\), then
\[
 F(t_n)^{3^{j+2}-t_n}\longrightarrow m^9.
\]
Its prime floor would eventually be either \(m^9-1\) or \(m^9\),
both composite. Thus every \(A^{3^j}\) is nonintegral. Continuity
of the floor at those points shows that \(A\in\mathcal W(3^k)\),
so \(A\ge M\). This proves \(\liminf F(t)\ge M\).

Write \(p_j=\floor{M^{3^j}}\). No \(M^{3^j}\) is integral:
otherwise the next prescribed power is composite. The closed
prime cells are strictly nested, since
\[
 p_j^3<p_{j+1},\qquad p_{j+1}+1<(p_j+1)^3.
\]
Equality at the second endpoint would give the composite integer
\((p_j+1)^3-1\). Given \(\epsilon>0\), choose \(n\) with
\(p_n\ge P_0\) and
\((p_n+1)^{1/3^n}<M+\epsilon\). Its closed cell lies strictly
inside all earlier prime cells. This finite collection of strict
containments persists after replacing \(3^j\) by \(3^j-t\), for
small enough \(|t|\). Continue from the same \(p_n\) using
Lemma~\ref{mi:continuation}. The resulting admissible base gives
\[
 F(t)\le(p_n+1)^{1/(3^n-t)}.
\]
Taking upper limits, then \(\epsilon\downarrow0\), proves
continuity.

The increasing change of variable \(A\mapsto A^d\) gives the
exact scaling identity \(\xi(dC)=\xi(C)^{1/d}\) for \(d>0\).
Thus the two normalized tails are \(F(3^{-h})\) and
\(F(-3^{-h})\). Finally,
\(F(t)^{3^j-t}\to M^{3^j}\), and the nonintegrality already
proved gives stability of any fixed finite set of floors.
\end{proof}

The unnormalized tail minima tend to \(1\).
Neither this convergence nor a common finite prime prefix
transfers transcendence to \(M\).

\section{Arithmetic of a hypothetically algebraic Mills tower}
\label{mi:cubic}

\subsection{The external reduction and the least-prime property}

\begin{theorem}[Saito's reduction: external input]
\label{mi:saito}
If \(M\) is algebraic, then for some integer \(m\ge1\),
\(\beta=M^{3^m}\) is a totally real cubic Pisot number. Its
conjugates can be written
\[
 \beta,\ -\rho,\ \tau,\qquad
 0<\sigma:=|\tau|<\rho<1,
 \qquad
 \rho\le\min(\sigma^{17/23},\beta^{-17/40}).
\]
\end{theorem}

This is the specialization of
\cite[Theorem~1.7]{SaitoVariants} used throughout this section.
All statements below about the actual Mills tower are implications
of \(M\in\A\), unless explicitly stated otherwise.
Put
\[
 p_r=\floor{\beta^{3^r}},\quad
 g_r=p_{r+1}-p_r^3,\quad x_r=p_r^3,\quad
 C=\Norm_{\Q(\beta)/\Q}(\beta),\quad N=|C|.
\]
On a sufficiently late tail,
\[
 p_r=\Tr(\beta^{3^r}),
\]
because \(-\rho^{3^r}+\tau^{3^r}\in(-1,0)\).

\begin{lemma}[Eventual least-prime choice]
\label{mi:least-prime}
For the actual least Mills constant, without assuming
algebraicity, each sufficiently late prime is the least prime
strictly above the cube of its predecessor. In the notation above,
\[
 p_{r+1}=\operatorname{nextprime}(p_r^3).
\]
\end{lemma}

\begin{proof}
Every sufficiently large prime has an infinite strict cubic
continuation by Lemma~\ref{mi:continuation}.
If a prime \(q\) satisfies \(p_r^3<q<p_{r+1}\), its inverse-power
cell is contained in that of \(p_r\), and all earlier classical
cells are preserved. Continue from \(q\).
The upper endpoint of the replaced cell is at most
\(p_{r+1}^{1/3^{m+r+1}}<M\).
The continued base is therefore smaller than \(M\) and satisfies
every original prime-floor condition, contradicting minimality.
\end{proof}

\subsection{Divisibility and the exact exponent}
\label{mi:gaps}

The following arithmetic consequences were recorded in
\cite{VMRC010}. Their elementary proofs are supplied here.

\begin{theorem}[Gap divisibility and transcendental exponent]
\label{mi:gap-theorem}
Under \(M\in\A\), the sufficiently late gap sequence satisfies
\[
 g_r\mid g_{r+1},\qquad
 \frac{9g_{r+1}}{g_r^3}\longrightarrow1,\qquad
 \frac{g_r}{p_r}\longrightarrow\infty.
\]
Moreover,
\[
 g_r\sim3x_r^\kappa,\qquad
 \kappa=\frac{2+\log\rho/\log\beta}{3},
\]
where \(\kappa\) is transcendental and
\[
 \frac12+\frac{\log N}{6\log\beta}
 <\kappa<\frac{21}{40}.
\]
\end{theorem}

\begin{proof}
At a given index, let \(X_1,X_2,X_3\) be the \(3^r\)-th powers
of the three conjugates, and write
\[
 p=X_1+X_2+X_3,\quad
 b=X_1X_2+X_1X_3+X_2X_3,\quad
 c=X_1X_2X_3.
\]
These are integers. Cubing the roots gives the Newton state map
\begin{equation}\label{mi:newton}
 p'=p^3-3bp+3c,\qquad
 b'=b^3-3cbp+3c^2,\qquad c'=c^3.
\end{equation}
If \(h=c-bp=g_r/3\), then
\[
 h=-(X_1+X_2)(X_1+X_3)(X_2+X_3).
\]
Thus
\[
 h'=h\,\mathcal F,\qquad
 \mathcal F=\prod_{i<j}(X_i^2-X_iX_j+X_j^2).
\]
The multiplier is a positive integer: it is a symmetric
polynomial in algebraic integers, and all three roots are real.
Hence \(g_r\mid g_{r+1}\).

Since \(|\tau|<\rho<1<\beta\), the same product identity yields
\[
 p_r\sim\beta^{3^r},\qquad
 h_r\sim(\beta^2\rho)^{3^r},\qquad
 g_r\sim3(\beta^2\rho)^{3^r}.
\]
This proves the cubic gap asymptotic and the formula for
\(\kappa\). The absolute norm satisfies
\[
 1\le N=\beta\rho\sigma<\beta\rho^2.
\]
Consequently
\(\rho>\sqrt{N/\beta}\), which gives the stated strict lower
bound for \(\kappa\), and also \(\beta\rho>1\).
The latter implies \(g_r/p_r\to\infty\).
Saito's bound initially gives \(\kappa\le21/40\).

It remains to determine the arithmetic nature of the exponent.
If \(t=\log\rho/\log\beta\) were rational, it would be negative,
so \(\beta^u\rho^v=1\) for some positive integers \(u,v\).
Equivalently, \(\beta^u(-\rho)^v=(-1)^v\).
The Galois group of an irreducible cubic contains a three-cycle;
choose the orientation sending \(\beta\) to \(-\rho\) and
\(-\rho\) to \(\tau\). Applying it and taking absolute values gives
\(\rho^u\sigma^v=1\), impossible since both factors are below \(1\).
Thus \(t\) is irrational.
If \(t\) were algebraic, the Gelfond--Schneider theorem would
make \(\beta^t=\rho\) transcendental, again a contradiction
\cite{Schneider1934}. Therefore \(t\), and hence \(\kappa\), is
transcendental. The rational endpoint \(21/40\) is excluded.
\end{proof}

In particular, \(g_r\mid g_{r+s}\) and
\[
 p_{r+s}\equiv p_r^{3^s}\pmod{g_r}\qquad(s\ge1).
\]
These necessary conditions are not mutually contradictory as
presently understood. The theorem is not an unconditional
transcendence proof for \(M\).

\subsection{A fixed three-term identity in the unit branches}
\label{mi:unit}

\begin{proposition}[Eliminating the second symmetric coordinate]
\label{mi:unit-polynomial}
Let \(p,q,r\) be the traces of the first, third and ninth powers
of three algebraic numbers whose elementary symmetric functions
are integers, and whose product is \(c\).
Then
\begin{align}
 P_c(p,q,r):={}&q(p^3+3c-q)^3
     +9p^3(r-q^3-3c^3) \notag\\
 &\hspace{13mm}-27cp^3q(p^3-q)=0.
 \label{mi:pc}
\end{align}
If the starting norm is \(c=\pm1\), the same polynomial applies
at every consecutive triple of sparse traces.
\end{proposition}

\begin{proof}
Let \(b\) be the initial second symmetric coordinate.
Equation~\eqref{mi:newton} gives
\[
 q=p^3-3bp+3c,\quad
 b'=b^3-3cbp+3c^2,\quad
 r=q^3-3b'q+3c^3.
\]
Substitution into \eqref{mi:pc} cancels all terms.
Odd cubing preserves \(c=\pm1\), proving the final statement.
\end{proof}

For a nonunit norm, the corresponding polynomial is
\(P_{C^{3^r}}\); its coefficients vary with \(r\).
Neither the unit assumption nor a fixed polynomial can therefore
be imposed on all hypothetical algebraic Mills candidates.
Even in the unit branches, transferring this identity to one of
the shifted sequences requires a further argument.

\section{Local restrictions and an exact nonreturn certificate}
\label{mi:local}

The next two results retain the hypotheses of the source
\cite{VMRC031}. Their scope is a specified unit branch, not all
possible algebraic Mills constants.

\subsection{Linear 3-adic depth}

\begin{theorem}[Exact depth at a unit fixed state]
\label{mi:depth}
Let \(\varepsilon\in\{1,-1\}\) and let
\[
 f(X)=X^3-aX^2+bX+\varepsilon
\]
be irreducible, with
\(a\equiv\varepsilon\pmod3\) and \(b\equiv-1\pmod3\).
Let \(T_r,B_r\) be the first and second elementary symmetric
functions of the \(3^r\)-th powers of its roots.
If \(s=v_3(a+\varepsilon b)\), then \(s\) is a finite positive
integer and
\[
 v_3(T_r-\varepsilon)=r+s\qquad(r\ge s).
\]
\end{theorem}

\begin{proof}
The norm is \(-\varepsilon\) and is unchanged by odd powering.
Put
\[
 x_r=T_r-\varepsilon,\quad y_r=B_r+1,\quad
 z_r=x_r+\varepsilon y_r,\quad w_r=x_r-\varepsilon y_r.
\]
Expansion of \eqref{mi:newton} gives
\begin{align*}
 z_{r+1}
 &=z_r(3+3\varepsilon w_r+x_r^2-\varepsilon x_ry_r+y_r^2),\\
 w_{r+1}
 &=w_r(9+3\varepsilon w_r+x_r^2+\varepsilon x_ry_r+y_r^2).
\end{align*}
Both \(x_r,y_r\) are divisible by \(3\) at every stage.
The first multiplier is \(3\pmod9\); the second is divisible by
\(9\). Since \(z_0=a+\varepsilon b\ne0\), we get
\[
 v_3(z_r)=r+s,\qquad v_3(w_r)\ge2r+1.
\]
Nonvanishing of \(z_0\) follows from
\(f(-\varepsilon)=-a-\varepsilon b\) and irreducibility.
For \(r\ge s\), the two valuations are different, and
\(2x_r=z_r+w_r\) proves the assertion.
\end{proof}

\subsection{A splitting restriction}

\begin{lemma}[Projective return]
\label{mi:return}
Suppose \(p\nmid c\operatorname{disc}(f)\), where \(c\) is the
norm of the roots, and \(p\mid T_r\).
In a finite splitting field modulo \(p\), let \(Y_i\) be the
roots raised to \(3^r\). If the multiplicative subgroup
generated by the ratios \(Y_i/Y_j\) has order prime to \(3\),
then \(p\mid T_{r+jL}\) for all \(j\ge0\), for some \(L\ge1\).
\end{lemma}

\begin{proof}
Choose \(L\) with \(3^L\equiv1\) modulo the exponent of the
ratio subgroup. At each \(L\)-step, the three powered roots
are multiplied by one common nonzero scalar. Their zero sum
therefore recurs.
\end{proof}

\begin{theorem}[Complete splitting in a specified norm-\(-1\) branch]
\label{mi:splitting}
Let \(\beta\) be a totally real cubic Pisot number whose small
conjugate of largest modulus is negative. Suppose its polynomial is
\[
 f(X)=X^3-aX^2+bX+1,\qquad
 a\equiv1,\ b\equiv-1\pmod3,\qquad v_3(a+b)=1.
\]
If \(\floor{\beta^{3^r}}\) is eventually prime, then \(f\)
splits completely modulo every sufficiently late such prime.
\end{theorem}

\begin{proof}
Eventually the floor is \(T_r\), and it exceeds every fixed
prime divisor of \(3\operatorname{disc}(f)\).
For a prime \(p=T_r\), Theorem~\ref{mi:depth} gives
\(p\equiv1\pmod3\) and \(v_3(p-1)=r+1\).

If \(f\) is irreducible modulo \(p\), its roots are
\(u,u^p,u^{p^2}\). Their ratios have orders dividing
\(p^2+p+1\), whose \(3\)-adic valuation is \(1\).
Raising to \(3^r\) removes this \(3\)-primary part.
Lemma~\ref{mi:return} then gives recurrent divisibility by \(p\).

If the factorization type is \(1+2\), the roots lie in
\(\mathbb F_{p^2}\), with one in \(\mathbb F_p\) and the other
two conjugate. Since \(v_3(p^2-1)=r+1\), the \(3\)-primary
parts of their \(3^r\)-th powers belong to \(\mu_3\).
The two conjugate roots have the same such part because
\(p\equiv1\pmod3\). The product of all three parts is \(1\),
by the norm \(-1\), so the remaining part is the same as well.
Every ratio consequently has order prime to \(3\), and
Lemma~\ref{mi:return} again applies.

In either nonsplit case, an arithmetic progression of later
traces is divisible by \(p\); these traces eventually exceed
\(p\), so the corresponding floors are composite. This
contradicts an eventually prime tail.
\end{proof}

The theorem does not exclude completely split prime traces.
In particular, it does not show that every candidate has a
recurrent prime divisor.

\subsection{A prime divisor that never returns}
\label{mi:nonreturn}

\begin{proposition}[A completely split nonreturning seed]
\label{mi:seed}
Let \(\beta\) be the largest root of
\[
 f(X)=X^3-211X^2-7X+1,
 \qquad T_r=\Tr(\beta^{3^r}).
\]
It is a totally real cubic Pisot number with conjugates
\(\beta,-\rho,\sigma\), where
\[
 0<\sigma<\rho<1,\qquad
 \rho\le\min(\sigma^{17/23},\beta^{-17/40}).
\]
Every \(T_r\) equals \(\floor{\beta^{3^r}}\).
Both \(T_0=211\) and \(T_1=p=9\,398\,359\) are prime, and \(f\)
splits completely modulo \(p\), but
\[
 p\nmid T_r\qquad(r\ge2).
\]
\end{proposition}

\begin{proof}
Sign changes locate roots in \((-1/10,0),(0,1),(211,212)\).
The rational-root test excludes \(\pm1\), so \(f\) is irreducible.
Its root sum and norm give
\[
 \rho-\sigma=\beta-211>0,\qquad \rho\sigma=\beta^{-1}.
\]
Moreover \(212^{17}<10^{40}\), whence
\(\rho<1/10<\beta^{-17/40}\).
Combining this with the norm identity gives
\(\rho^{23}<\sigma^{17}\).
For every positive odd integer \(n\),
\(-\rho^n+\sigma^n\in(-1,0)\), which proves the floor identity.

Newton's formulas give \(T_1=p\) and \(B_1=-4771\).
Trial division to the respective integer square roots certifies
the two primality assertions. The three distinct roots modulo
\(p\) are
\[
 -1\,584\,786,\qquad -1\,942\,624,\qquad 3\,527\,621.
\]
For norm \(-1\), the modular state map is
\[
 \Phi(t,b)=(t^3-3tb-3,\ b^3+3tb+3)\pmod p.
\]
Starting at the index-\(1\) state \((0,-4771\bmod p)\), its
next state is
\[
 S_1=(9\,398\,356,\ 7\,830\,596).
\]
The integer verifier reproduced in
Appendix~\ref{mi:app:nonreturn} checks that the orbit from \(S_1\)
has exactly \(85\,712\) distinct states, returns to \(S_1\), and
has no zero first coordinate. This finite exhaustive certificate
proves nondivisibility at every index \(r\ge2\).
\end{proof}

This is a counterexample to universal return of a chosen prime
seed, not to the existence of composite terms. Indeed,
\[
 T_2=830\,149\,079\,789\,343\,894\,643
     =7\cdot31\cdot3\,825\,571\,796\,264\,257\,579.
\]
The example also satisfies \(a\equiv1\), \(b\equiv-1\pmod3\)
and \(v_3(a+b)=1\). It shows why the local restrictions above
do not by themselves force return of an arbitrarily selected
prime trace.

\section{Quantitative limits of the transfer to the original constant}
\label{mi:limits}

\subsection{Keeping three prime outputs versus forcing polynomial escape}
\label{mi:small-shift}

This subsection analyzes the particular perturbation used in the
discussion. It is not an impossibility theorem for all perturbation
methods.

\begin{lemma}[Necessary bounds for a preserved triple]
\label{mi:preserve}
Work on the hypothetical algebraic Mills tail of
Theorem~\ref{mi:gap-theorem}, and put
\(a=-\log\rho/\log\beta\in(17/40,1/2)\).
Let \(B\) stay in a fixed compact subinterval of \((1,\infty)\).
Suppose the clock \(3^k-t\), at three consecutive indices
\(k,k+1,k+2\), has the same outputs \(p,q,r\) as the
unshifted tail, with \(p\to\infty\). Then
\[
 t>0\ \Longrightarrow\ t=O(p^{-3(1+a)}),\qquad
 t<0\ \Longrightarrow\ |t|=O(p^{-3}).
\]
In particular, \(|t|p^3=O(1)\).
\end{lemma}

\begin{proof}
For any two preserved outputs \(p,q\), the exponent identity
\((3^{k+1}-t)-3(3^k-t)=2t\) gives the exact inequalities
\begin{equation}\label{mi:ratio-window}
 \frac{q}{(p+1)^3}<B^{2t}<\frac{q+1}{p^3}.
\end{equation}
Since \(q-p^3\asymp p^{2-a}\), the upper bound implies
\(t=O(p^{-1-a})\) when \(t>0\).
The lower bound, whose logarithm is
\(-3/p+O(p^{-1-a}+p^{-2})\), implies
\(|t|=O(p^{-1})\) when \(t<0\).
Apply the same estimates to the pair \(q,r\), using
\(q\sim p^3\) and the same fixed \(a\).
\end{proof}

For comparison, the fixed polynomial \eqref{mi:pc} has a
directly verifiable quantitative escape estimate.

\begin{proposition}[One sufficient escape scale for \(P_c\)]
\label{mi:escape-scale}
Fix \(c=\pm1\), let \(s\to\infty\), \(L\to1\), and put
\[
 u=\floor{s},\qquad v=\floor{Ls^3},\qquad
 w=\floor{L^4s^9}.
\]
If \(|L-1|s^3\to\infty\), then \(P_c(u,v,w)\ne0\)
eventually.
\end{proposition}

\begin{proof}
Exact substitution gives
\begin{align*}
 P_c(s,Ls^3,L^4s^9)
 =(L-1)\bigl\{&
 s^{12}L(4L-1)(2L+1)
 +9cs^9L(L+2)\\
 &-27c^2s^6L+27c^3s^3\bigr\}.
\end{align*}
The leading bracket is asymptotic to \(9s^{12}\).
Writing the weighted-degree-\(12\) part of \(P_c\) as
\[
 H(x,y,z)=y(x^3-y)^3+9x^3(z-y^3),
\]
one checks \(H_x(s,s^3,s^9)=H_{xx}(s,s^3,s^9)=0\).
Taylor expansion in the first coordinate, and first-order
bounds in the other two, therefore give the uniform floor-error
estimate
\[
 P_c(u,v,w)-P_c(s,Ls^3,L^4s^9)
 =O\bigl(|L-1|s^{11}+s^9\bigr).
\]
The terms of lower weighted degree contribute only \(O(s^9)\).
Under the stated hypothesis this error is
\(o(|L-1|s^{12})\), proving nonvanishing.
\end{proof}

For a perturbed clock, \(s=B^{3^k-t}\) and \(L=B^{2t}\).
When \(B\) remains in the stated compact set and \(t\to0\),
\(|L-1|\asymp|t|\). Thus the sufficient condition just proved
requires \(|t|p^3\to\infty\), whereas preservation of the three
outputs gives \(|t|p^3=O(1)\).
These two particular estimates do not overlap.
Pointwise fixed-base avoidance cannot be silently made uniform
when the base and the shift vary with the index.

\subsection{Intermediate prime conditions remove a skipped-level gain}
\label{mi:corridor}

\begin{lemma}[The ancestor of an integer root cell]
\label{mi:ancestor}
For integers \(q,D\ge1\), the integer part is constant on
\([q^{1/D},(q+1)^{1/D})\), with value
\(\floor{q^{1/D}}\).
\end{lemma}

\begin{proof}
If \(j=\floor{q^{1/D}}\), then
\(j^D\le q<(j+1)^D\). Since both \(q\) and \((j+1)^D\)
are integers, \(q+1\le(j+1)^D\); the half-open endpoint proves
the claim.
\end{proof}

\begin{proposition}[The terminal search interval]
\label{mi:terminal}
On a sufficiently late least-Mills tail, fix a depth \(d\ge2\).
Among prime terminal labels \(q<p_{r+d}\), preservation of the first
\(j\) descendants leaves the geometric search interval
\[
 I_{r,d}^{(j)}
   =(p_{r+j}^{\,3^{d-j}},\,p_{r+d}),
 \qquad 0\le j<d.
\]
For \(j<d-1\), requiring the next ancestor to be prime restricts
this interval to \(I_{r,d}^{(j+1)}\).
After all intermediate conditions are imposed, the remaining
interval is
\[
 I_{r,d}^{(d-1)}=(p_{r+d-1}^3,\,p_{r+d}).
\]
This statement does not assert the existence of a prime in any
of these intervals.
\end{proposition}

\begin{proof}
By Lemma~\ref{mi:ancestor}, a terminal label determines each
earlier integer floor exactly. Below
\(p_{r+j+1}^{3^{d-j-1}}\), its next ancestor lies between
\(p_{r+j}^3\) and \(p_{r+j+1}-1\).
Every integer in this range is composite, by the least-prime
property, including the initial cube.
Above this threshold and below \(p_{r+d}\), the next ancestor
is exactly \(p_{r+j+1}\): the upper bound follows from the
nesting inequality
\(p_{r+d}<(p_{r+j+1}+1)^{3^{d-j-1}}\).
Iteration proves the claim.
\end{proof}

The size of the discarded portion can be quantified.
With \(a=-\log\rho/\log\beta\), \(X=\beta^{3^{r+d}}\), and
fixed \(d,j\), binomial expansion gives
\[
 |I_{r,d}^{(j)}|
 \sim3^{d-j}X^{\,1-(1+a)/3^{d-j}}.
\]
For example, a raw two-level exponent lies in
\((5/6,101/120)\), whereas the final exponent is
\(\kappa\in(1/2,21/40)\).
More precisely,
\[
 \frac{|I_{r,d}^{(d-1)}|}{|I_{r,d}^{(0)}|}
 \sim3^{1-d}
 \left(\frac{\rho}{\beta}\right)^{(3^{d-1}-1)3^r}
 \longrightarrow0.
\]
To verify the asymptotic, write
\(p_{r+j}=Y-e_{r+j}\), \(Y=\beta^{3^{r+j}}\),
where \(e_{r+j}\sim\rho^{3^{r+j}}\).
The leading term of
\(Y^{3^{d-j}}-e_{r+d}-(Y-e_{r+j})^{3^{d-j}}\)
is \(3^{d-j}Y^{3^{d-j}-1}e_{r+j}\); the other terms are smaller.
Thus the apparent large-interval improvement is removed by
the intermediate prime constraints \cite{VMRC075}.

A collection of many intervals would require an additional
quantitative statement about their number and distribution.
No such statement sufficient for the present sparse orbit was
established in the discussion.

\subsection{The analytic input that would finish this route}
\label{mi:rh}

Let \(N(\sigma,T)\) count zeta zeros with real part at least
\(\sigma\) and imaginary part in \((0,T]\).
The density hypothesis is the estimate
\[
 N(\sigma,T)\ll_\epsilon T^{2(1-\sigma)+\epsilon}
 \qquad(1/2\le\sigma\le1).
\]
Saito proves that DH, and hence RH, implies the transcendence
of \(M\) \cite[Theorem~1.8]{SaitoVariants}.
The role of the short-interval input is transparent here:
an available prime between \(x\) and \(x+x^\theta\), for every
fixed \(\theta>1/2\) and every sufficiently large \(x\), can be
applied with \(1/2<\theta<\kappa\).
It then places a prime strictly between \(p_r^3\) and
\(p_{r+1}\), contradicting Lemma~\ref{mi:least-prime}.

The OpenAI scope document discussed in the conversation states
a zero-free half-plane \(\Re s>7/8\) for the zeta function and
the listed \(L\)-functions \cite{OpenAIQRH}.
That scope statement does not supply the density estimate
displayed above, or a prime-insertion theorem at every exponent
\(\theta>1/2\).
No deduction of the required input from that statement is
provided here. It is therefore not used as a replacement for
DH or RH.

The two unresolved transfers are explicit: a shifted
polynomial-avoidance theorem has not been transferred to the
actual unshifted minimum, and the necessary prime in its
hypothetical algebraic gap has not been forced unconditionally.

\section{Reciprocal series and an irrationality transfer criterion}
\label{rs:section}

This section supplies the direct connection between the two themes.  It
covers Mills item M2 and gamma item G6.  The conclusions concern reciprocal
series built from shifted power floors; an identity expressing Euler's
constant by one of these series is a separate, unproved requirement.

\subsection{Four clocks and an integer-tail asymptotic}

Fix $A>1$, and consider one of the four exponent sequences
\begin{equation}\label{rs:clocks}
 C_k=3^k+\varepsilon\quad\hbox{or}\quad C_k=3^k+\varepsilon k,
 \qquad \varepsilon\in\{-1,1\},\quad k\ge1.
\end{equation}
Put $p_k=\floor{A^{C_k}}$.  For the intended applications these integers
are primes.  In the proofs it is enough that $p_k\ge1$; they are at least
$2$ for all sufficiently large $k$.  Define
\[
 S_1=\sum_{k\ge1}\frac1{p_k},\qquad
 S_2=\sum_{k\ge1}\frac1{p_k^2},\qquad
 P_n=\prod_{k=1}^n p_k.
\]
Both series converge absolutely.  Finite initial terms make no difference
to linear independence together with $1$.

\begin{lemma}[Normalized product and tail]\label{rs:tail}
The product
\[
 \Omega=\prod_{k\ge1}\left(\frac{p_k}{A^{C_k}}\right)^2
\]
converges to a positive real number.  If $a,b$ are fixed real numbers
with $a\ne0$, then, for some $c_0>0$,
\begin{equation}\label{rs:integer-asymptotic}
 P_n^2\sum_{k>n}\left(\frac a{p_k}+\frac b{p_k^2}\right)
 =a\Omega A^{d_n}+O\bigl(\exp(-c_0 3^n)\bigr),
\end{equation}
where
\begin{equation}\label{rs:dn}
 d_n=2\sum_{k=1}^n C_k-C_{n+1}
 =\begin{cases}
 -3+\varepsilon(2n-1),&C_k=3^k+\varepsilon,\\
 -3+\varepsilon(n^2-1),&C_k=3^k+\varepsilon k.
 \end{cases}
\end{equation}
The implied constants may depend on $A,a,b$ and the chosen clock.
\end{lemma}

\begin{proof}
The relative error in a floor is
$p_k/A^{C_k}=1+O(A^{-C_k})$.  The sum of these errors converges, and no
factor is zero.  The usual estimate for a convergent infinite product gives
\[
 P_n^2=\Omega A^{2\sum_{k\le n}C_k}
       \bigl(1+O(A^{-C_{n+1}})\bigr).
\]
The first tail term satisfies
$a/p_{n+1}=aA^{-C_{n+1}}(1+O(A^{-C_{n+1}}))$.
The $b$-term at that index is $O(A^{-2C_{n+1}})$, and the remaining
tail is $O(A^{-C_{n+2}})$.  Since $C_{n+2}/C_{n+1}\to3$, these estimates
yield the claimed main term with error $O(A^{d_n-C_{n+1}})$.
Formula \eqref{rs:dn} follows by summing the exponents.  Because
$d_n=O(n^2)$ and $C_{n+1}$ grows as $3^{n+1}$, the error is bounded
by $O(\exp(-c_0 3^n))$ after decreasing $c_0$ if necessary.
\end{proof}

\subsection{Linear independence of two reciprocal moments}

\begin{theorem}[Two reciprocal series]\label{rs:independence}
For either negative clock $3^k-1$ or $3^k-k$, and every $A>1$,
\[
 \dim_{\Q}\langle 1,S_1,S_2\rangle=3.
\]
For either positive clock $3^k+1$ or $3^k+k$, the same conclusion holds
provided $A^2\notin\Q$.  In particular it holds for each of the four
shifted prime-generating bases when that base is transcendental.
\end{theorem}

\begin{proof}
Suppose $c+aS_1+bS_2=0$ with integers $a,b,c$.  Then
\begin{align*}
 I_n&=P_n^2\sum_{k>n}\left(\frac a{p_k}+\frac b{p_k^2}\right)\\
 &=-cP_n^2-\sum_{k\le n}
 \left(a\frac{P_n^2}{p_k}+b\frac{P_n^2}{p_k^2}\right)\in\Z.
\end{align*}
If $a=0$ and $b\ne0$, these are nonzero integers, while
\[
 |I_n|\ll A^{2\sum_{k\le n}C_k-2C_{n+1}}
       =A^{d_n-C_{n+1}}\longrightarrow0.
\]
This is impossible.  It remains to consider $a\ne0$.  For all sufficiently
large $k$, the sign of $a/p_k+b/p_k^2$ equals the sign of $a$, so $I_n\ne0$
eventually.  Lemma~\ref{rs:tail} applies.

For the two negative clocks, $d_n\to-\infty$, and again
$0<|I_n|\to0$ contradicts integrality.

For $C_k=3^k+1$, we have $d_n=2n-4$.  The integer
\[
 J_n=I_nI_{n+2}-I_{n+1}^2
\]
tends to zero: its leading terms cancel, and every remaining term is
$O(\exp(O(n)-c_0 3^n))$.  Thus $J_n=0$ for all sufficiently large $n$.
Since $I_n\ne0$, the rational ratios $I_{n+1}/I_n$ are then constant.
Their limit, computed from \eqref{rs:integer-asymptotic}, is $A^2$.
This contradicts $A^2\notin\Q$.

For $C_k=3^k+k$, we have $d_n=n^2-4$.  Instead use the integer
\[
 J_n=I_nI_{n+2}^3-I_{n+1}^3I_{n+3}.
\]
The identity $d_n+3d_{n+2}=3d_{n+1}+d_{n+3}$ cancels the leading terms.
The remaining error is $O(\exp(O(n^2)-c_0 3^n))$, which tends to zero.
Therefore $J_n=0$ eventually.  The nonzero rational quantities
\[
 R_n=\frac{I_nI_{n+2}}{I_{n+1}^2}
\]
then satisfy $R_{n+1}=R_n$, whereas their limit is
$A^{d_n+d_{n+2}-2d_{n+1}}=A^2$.  This is the same contradiction.
Consequently $a=b=0$, and the original relation forces $c=0$.
\end{proof}

\begin{corollary}\label{rs:irrational-combinations}
Under the hypotheses of Theorem~\ref{rs:independence}, every nonzero
rational linear combination $aS_1+bS_2$ is irrational.
\end{corollary}
\begin{proof}
A rational value would give a nontrivial rational relation among
$1,S_1,S_2$.
\end{proof}

\begin{remark}[Scope and the role of primality]\label{rs:scope}
The theorem applies to each shifted sequence separately.  It does not
prove joint independence across the four different sequences, algebraic
independence of $S_1,S_2$, or transcendence of either series.  Primality is
not needed in the proof: the decisive inputs are the floor approximation
and the exponent clock.  The positive cases need only $A^2\notin\Q$,
which is weaker than the transcendence hypothesis used in the original
discussion.  No assertion of historical priority is made for this
integer-tail argument.
\end{remark}

\subsection{A sparse expansion criterion for irrationality}

\begin{theorem}[Bounded integer coefficients on negative clocks]
\label{rs:transfer}
Let $C_k=3^k-1$ or $C_k=3^k-k$, let $A>1$, and put
$p_k=\floor{A^{C_k}}$.  Suppose $r\in\Q$ and $a_k\in\Z$ satisfy
\[
 |a_k|\le\sqrt{p_k},\qquad
 a_k\ne0\text{ for infinitely many }k.
\]
Then the absolutely convergent number
\[
 x=r+\sum_{k\ge1}\frac{a_k}{p_k}
\]
is irrational.
\end{theorem}

\begin{proof}
Suppose $x\in\Q$, and choose an integer $b>0$ clearing the denominators
of both $x$ and $r$.  Then
\[
 J_n=bP_n\sum_{k>n}\frac{a_k}{p_k}
   =bP_n(x-r)-b\sum_{k\le n}a_k\frac{P_n}{p_k}\in\Z.
\]
Because $C_{k+1}/C_k\to3$,
\[
 |J_n|\le bP_n\sum_{k>n}p_k^{-1/2}
 \ll A^{\sum_{k\le n}C_k-C_{n+1}/2}
 =A^{d_n/2}\longrightarrow0.
\]
For every sufficiently large $j$ with $a_j\ne0$, integrality of $a_j$
gives $|a_j/p_j|\ge p_j^{-1}$, while
\[
 \sum_{k>j}\left|\frac{a_k}{p_k}\right|
 \le\sum_{k>j}p_k^{-1/2}
 \ll p_{j+1}^{-1/2}=o(p_j^{-1}).
\]
Hence the tail beginning at $j$ is nonzero, so $J_{j-1}\ne0$ for
infinitely many $j$.  This contradicts $J_n\in\Z$ and $J_n\to0$.
\end{proof}

\begin{corollary}[The precise missing gamma identity]\label{rs:gamma-link}
If Euler's constant admits an identity
\[
 \gamma=r+\sum_{k\ge1}\frac{a_k}{p_k}
\]
satisfying all hypotheses of Theorem~\ref{rs:transfer}, then $\gamma$
is irrational.
\end{corollary}

The theorem proves the implication; the required representation of
$\gamma$ has not been constructed here.  A formal expansion with arbitrary
real coefficients would not suffice.  Neither would an expansion whose
integer coefficients violate the stated square-root bound.  These are
the exact requirements needed to preserve integrality, smallness, and
nonvanishing simultaneously.

\subsection{What the four transcendence inputs do not imply}

Individual transcendence of four constants does not imply their joint
algebraic independence: for a transcendental $t$, the four numbers
$t,t+1,t^2,t^3$ are all transcendental and are algebraically dependent.
Likewise, the arithmetic property of a limit does not follow from that
property of every term: $1+\pi/n$ is transcendental for every $n$ but
converges to $1$.  These elementary observations explain why an exact
transfer such as Corollary~\ref{rs:gamma-link}, or the quantified
structural statements in the other sections, is necessary.

\section{Primitive arithmetic of logarithm-free approximants}
\label{ga:section}

This section contains the approximation and denominator results denoted
G7--G10 in the result register.  It separates three different assertions:
an exact positive error formula, a theorem excluding a specified
asymptotic range of degrees, and a finite exact computation including degrees not controlled by that
asymptotic consequence.  None of these assertions proves that $\gamma$ is rational or
irrational.

\subsection{The arithmetic target}
\label{ga:target}

\begin{proposition}[A sufficient irrationality criterion]
\label{ga:irrationality}
Suppose that $P_j\in\Z$, $Q_j\in\N$, and
\[
             0<|Q_j\gamma-P_j|\longrightarrow0.
\]
Then $\gamma\notin\Q$.
\end{proposition}
\begin{proof}
If $\gamma=a/b$ with $a\in\Z$ and $b\in\N$, then
$b(Q_j\gamma-P_j)=aQ_j-bP_j$ is a nonzero integer tending to zero.
\end{proof}

Thus convergence of $P_j/Q_j$ to $\gamma$ is not by itself the
arithmetic target.  The error must still tend to zero after multiplication
by an integer denominator.  Taking the denominator in lowest terms gives
the most favorable normalization for that particular rational
approximant.

\subsection{Binet's measure and its Gaussian approximants}
\label{ga:binet}

For $s>0$, put
\begin{equation}\label{ga:measure}
 d\mu(u)=\frac{du}{e^{2\pi\sqrt u}-1},\qquad
 S(s)=\int_0^\infty\frac{d\mu(u)}{s+u}.
\end{equation}
Every moment is finite and rational:
\begin{equation}\label{ga:moments}
 \mu_j:=\int_0^\infty u^j\,d\mu(u)
 =\frac{2\Gamma(2j+2)\zeta(2j+2)}{(2\pi)^{2j+2}}
 =\frac{|B_{2j+2}|}{2j+2}\qquad(j\ge0).
\end{equation}
Indeed, substitute $u=t^2$, expand $(e^{2\pi t}-1)^{-1}$ into
its positive geometric series, and integrate termwise.  Euler's
formula for the even zeta values gives the last equality.
The first four moments are
$1/12,1/120,1/252,1/240$.

The differentiated Binet formula
\cite[Eq.~5.9.15]{NISTDLMF}, together with
$\psi(n)=H_n-\gamma-1/n$, gives
\begin{equation}\label{ga:binet-identity}
 H_n-\frac1{2n}+S(n^2)=\gamma+\log n
 \qquad(n\in\N).
\end{equation}

Let $p_m\in\Q[u]$ be the monic orthogonal polynomial of degree $m$
for $\mu$, and set $p_0=1$.  The moment matrices are positive
definite, so these polynomials exist uniquely.  All their zeros are
simple and belong to $(0,\infty)$; this follows from the usual
sign-change argument for orthogonality with a strictly positive
measure.  Define
\begin{equation}\label{ga:R-definition}
 R_m(s)=\frac1{p_m(-s)}
 \int_0^\infty
       \frac{p_m(-s)-p_m(u)}{s+u}\,d\mu(u),\qquad R_0=0.
\end{equation}
Polynomial division shows that $R_m\in\Q(s)$, with numerator
degree at most $m-1$ and denominator degree $m$.  It is the
$[m-1/m]$ Gaussian Pad\'e approximant of the formal moment expansion
$S(s)\sim\sum_{j\ge0}(-1)^j\mu_js^{-j-1}$ at infinity.
No convergence of that formal expansion is assumed.

\begin{proposition}[Positive error and a rational lower certificate]
\label{ga:positive-error}
For $m\ge0$ and $s>0$,
\begin{equation}\label{ga:error}
 E_m(s):=S(s)-R_m(s)
 =\frac1{p_m(-s)^2}
       \int_0^\infty\frac{p_m(u)^2}{s+u}\,d\mu(u)>0.
\end{equation}
The function $E_m$ is strictly decreasing on $(0,\infty)$.  Moreover,
\begin{equation}\label{ga:consecutive}
 \begin{gathered}
 0<R_{m+1}(s)-R_m(s)
 =\frac{h_m}{V_m(s)V_{m+1}(s)}<E_m(s),\\
 h_m=\int p_m^2\,d\mu,\qquad V_m(s)=(-1)^mp_m(-s)>0.
 \end{gathered}
\end{equation}
\end{proposition}
\begin{proof}
Subtract \eqref{ga:R-definition} from $S(s)$.  Orthogonality gives
\[
 \int p_m(u)
       \frac{p_m(-s)-p_m(u)}{s+u}\,d\mu(u)=0,
\]
because the quotient is a polynomial in $u$ of degree $m-1$.
This proves \eqref{ga:error}.  Its numerator integral strictly
decreases with $s$, while $p_m(-s)^2$ increases for $m\ge1$;
for $m=0$ the numerator alone is strictly decreasing.

For completeness, write the monic three-term recurrence as
\[
 p_{m+1}(u)=(u-a_m)p_m(u)-b_mp_{m-1}(u),\qquad
 a_m=\frac{\int up_m^2\,d\mu}{h_m},\quad
 b_m=\frac{h_m}{h_{m-1}}>0\quad(m\ge1).
\]
The denominator polynomials $V_m$ and the corresponding numerators
$U_m$ of $R_m=U_m/V_m$ satisfy the same recurrence
$X_{m+1}=(s+a_m)X_m-b_mX_{m-1}$, with
\[
 V_0=1,\quad V_1=s+a_0,\qquad U_0=0,\quad U_1=\mu_0.
\]
Consequently
$U_{m+1}V_m-U_mV_{m+1}=h_m$ by induction.  This gives the
middle identity in \eqref{ga:consecutive}; the last inequality follows
from $E_m=R_{m+1}-R_m+E_{m+1}$.
\end{proof}

\subsection{An exact two-scale logarithm cancellation}
\label{ga:two-scale}

Let $N\ge1$ and $n=2^N$.  Define
\begin{equation}\label{ga:r}
\begin{split}
 r_{m,N}={}&(N+1)
   \left(H_n-\frac1{2n}+R_m(n^2)\right)\\
 &-N\left(H_{2n}-\frac1{4n}+R_m(4n^2)\right).
\end{split}
\end{equation}
This construction, used in \cite{VMRC019}, makes the logarithms in
\eqref{ga:binet-identity} cancel exactly:
$(N+1)\log n-N\log(2n)=0$.

\begin{theorem}[Logarithm-free positive rational approximants]
\label{ga:logfree}
For every $m,N\ge1$, one has $r_{m,N}\in\Q$ and
\begin{equation}\label{ga:gamma-error}
\begin{split}
 \gamma-r_{m,N}
 &=(N+1)E_m(n^2)-NE_m(4n^2)\\
 &=E_m(n^2)+N\{E_m(n^2)-E_m(4n^2)\}\\
 &>E_m(n^2)>R_{m+1}(n^2)-R_m(n^2)>0.
\end{split}
\end{equation}
In particular, if $r_{m,N}=P_{m,N}/Q_{m,N}$ is in lowest terms
with $Q_{m,N}>0$, then
\begin{equation}\label{ga:exact-lower}
 Q_{m,N}\gamma-P_{m,N}
 >
 L_{m,N}:=
 Q_{m,N}\{R_{m+1}(n^2)-R_m(n^2)\}>0,
\end{equation}
and $L_{m,N}$ is a rational number that can be compared exactly
with $1$.
\end{theorem}
\begin{proof}
Substitute \eqref{ga:binet-identity} into \eqref{ga:r}, then apply
Proposition~\ref{ga:positive-error}.  Rationality follows from the
rationality of all the moments.
\end{proof}

The approximants do converge to $\gamma$ as $N\to\infty$, even if
$m=m(N)$ varies arbitrarily in $\N$, since
\[
 0<\gamma-r_{m,N}
 <(N+1)S(n^2)\le\frac{(N+1)\mu_0}{n^2}.
\]
The five requirements in the research program therefore have the
following exact status.

\begin{center}
\begin{tabular}{@{}p{0.46\linewidth}p{0.46\linewidth}@{}}
\hline
Requirement & Status for \eqref{ga:r}\\
\hline
Exact removal of logarithms & Established.\\
Strictly nonzero error & Established by positivity.\\
Error tends to zero after multiplication by the reduced denominator
 & Not established; see the results below.\\
Infinitely many approximants & Established.\\
No unproved auxiliary hypothesis & Established for the construction.\\
\hline
\end{tabular}
\end{center}

\subsection{A denominator barrier for slowly increasing degrees}
\label{ga:slow-degree}

Write $\vartheta(x)=\sum_{\ell\le x,\ \ell\ {\rm prime}}\log\ell$.
The following statement is the slow-degree result in
\cite{VMRC019}; we give its denominator and error arguments explicitly.

\begin{theorem}[Slow-degree Stieltjes barrier]
\label{ga:slow-barrier}
There is an absolute constant $C>0$ such that, for every
$m,N\ge1$ and $n=2^N$,
\begin{equation}\label{ga:barrier}
\begin{split}
 \log\!\left\{Q_{m,N}(\gamma-r_{m,N})\right\}
 \ge{}&\vartheta(2n)-\vartheta(n)\\
 &-C\{m^2\log(m+1)+(m+1)\log n\}.
\end{split}
\end{equation}
Consequently, as $N\to\infty$,
\[
 m^2\log(m+1)+(m+1)\log n=o(n)
 \quad\Longrightarrow\quad
 Q_{m,N}(\gamma-r_{m,N})\longrightarrow+\infty.
 \tag{*}
\]
In particular, $(\ast)$ holds if
$m=o(\sqrt{n/\log n})$.
\end{theorem}
\begin{proof}
We first control the rational coefficients independently of $n$.
The von Staudt--Clausen denominator theorem shows that a common
denominator $\mathfrak d_m$ of $\mu_0,\ldots,\mu_{2m}$ can be chosen
to divide
\[
 \operatorname{lcm}(1,\ldots,4m+2)
       \prod_{\ell\le4m+3,\ \ell\ {\rm prime}}\ell.
\]
Thus $\log\mathfrak d_m=O(m)$.  The factorial expression in
\eqref{ga:moments} implies that the integerized moments have
logarithmic size $O(m\log(m+1))$.
Apply Cramer's rule to the $m$ equations
$\int u^jp_m(u)\,d\mu(u)=0$, $0\le j<m$.
After multiplication of the moment matrix by $\mathfrak d_m$,
its determinant is a nonzero integer, and every determinant involved
is bounded by $m!$ times an $m$th power of the largest entry.
It follows that the rational coefficients of $p_m$, and their
numerators and denominators, have logarithmic height
$O(m^2\log(m+1))$.

In particular, there are integer polynomials
$\widetilde U_m,\widetilde V_m$ of coefficient logarithmic height
$O(m^2\log(m+1))$, with $\deg\widetilde V_m=m$, such that
\[
 R_m(s)=\frac{\widetilde U_m(s)}
                {\mathfrak d_m\widetilde V_m(s)}.
\]
The polynomial $\widetilde V_m$ is an integer multiple of
$p_m(-s)$, so it has no positive real zero.  Its evaluation satisfies
\[
 \log|\widetilde V_m(t)|
 \le C_1m^2\log(m+1)+m\log(1+t)\qquad(t>0).
\]

Now let $\ell$ be a prime with $n<\ell\le2n$.
Since $n$ is a power of $2$, the rational terms $1/(2n)$ and $1/(4n)$
are integral at $\ell$.  The harmonic term $H_n$ is also integral
there, whereas $H_{2n}$ has exactly one summand with denominator
divisible by $\ell$, namely $1/\ell$.  Moreover, $\ell>N$.
Unless
\[
 \ell\mid
 \mathfrak d_m\widetilde V_m(n^2)\widetilde V_m(4n^2),
\]
both Pad\'e evaluations are integral at $\ell$, and \eqref{ga:r}
has $\ell$-adic valuation $-1$, coming from $-N/\ell$.
Every such nonexceptional prime therefore divides its actual
reduced denominator $Q_{m,N}$.  The total logarithm of the
exceptional primes is at most the logarithm of the displayed
nonzero integer.  Hence
\begin{equation}\label{ga:denominator-lower}
 \log Q_{m,N}\ge\vartheta(2n)-\vartheta(n)
 -C_2\{m^2\log(m+1)+m\log n\}.
\end{equation}

It remains to give a lower bound on the positive error.  The density
of $\mu$ on $[1,2]$ is at least
$c_\mu=(e^{2\pi\sqrt2}-1)^{-1}>0$.
Among monic polynomials of degree $m$ on an interval of length $1$,
the shifted monic Legendre polynomial minimizes the squared
$L^2$ norm.  Consequently
\[
 \int_1^2p_m(u)^2\,du
 \ge\frac1{(2m+1)\binom{2m}{m}^{\,2}}
 \ge\frac1{(2m+1)16^m}.
\]
Formula \eqref{ga:error} now yields
\[
 E_m(n^2)\ge
 \frac{c_\mu}
 {(n^2+2)(2m+1)16^m\,p_m(-n^2)^2}.
\]
The coefficient height bound already proved implies
\begin{equation}\label{ga:error-lower}
 \log E_m(n^2)\ge
 -C_3\{m^2\log(m+1)+(m+1)\log n\}.
\end{equation}
Combine \eqref{ga:denominator-lower},
\eqref{ga:error-lower}, and \eqref{ga:gamma-error}.
Finally the prime number theorem gives
$\vartheta(2n)-\vartheta(n)=n+o(n)$, proving $(\ast)$.
\end{proof}

This is an asymptotic exclusion for the stated degree range and this
particular two-scale family.  It makes no assertion that every high
degree, different linear combination, or different approximation
method must fail.

\subsection{Exact computation beyond the slow-degree range}
\label{ga:finite-audit}

To inspect degrees not covered by the preceding asymptotic
consequence, the calculation used
\[
 n\in\{2,4,8,16,32,64\},\qquad 1\le m\le64.
\]
The monic orthogonal recurrence was evaluated with exact rational
arithmetic.  For each of these $384$ pairs, \eqref{ga:r} was reduced
to lowest terms before computing $L_{m,N}$ in
\eqref{ga:exact-lower}.

\begin{proposition}[Finite exact certificate]
\label{ga:384}
For all $384$ pairs in the displayed range,
\[
                 L_{m,N}>1,
 \qquad\text{and therefore}\qquad
                 Q_{m,N}\gamma-P_{m,N}>1.
\]
\end{proposition}
\begin{proof}[Computational proof]
The standard-library verifier included in the computational appendix
constructs the Bernoulli numbers as rational numbers, computes the
monic orthogonal recurrence, and compares the numerator and
denominator of each positive rational $L_{m,N}$.
The test is an integer comparison; no numerical value of $\gamma$
is an input.  The archived output contains exactly $384$ successful
comparisons.  The small cases $m=1,2,3$, $n=2,4$ were additionally
checked by solving the moment orthogonality equations directly,
independently of that recurrence.
\end{proof}

For an explicit certificate visible without running a large
calculation, the first recurrence coefficients are
\[
 a_0=\frac1{10},\qquad
 a_1=\frac{871}{790},\qquad b_1=\frac{79}{2100}.
\]
Thus
\[
 R_1(s)=\frac5{6(10s+1)},\qquad
 R_2(s)=
 \frac{(s+871/790)/12}
 {(s+1/10)(s+871/790)-79/2100}.
\]
At $n=4$, $N=2$, $m=1$ these give
\begin{equation}\label{ga:small-certificate}
 r_{1,2}=\frac{3574141}{6192060},\qquad
 L_{1,2}=\frac{4000481}{913489}>4,
\end{equation}
and hence
$6192060\gamma-3574141>4000481/913489>4$.

\begin{center}
\begin{tabular}{@{}rrrr@{}}
\hline
$n$ & $m$ & Digits of the reduced $Q_{m,N}$
 & $\log_{10}L_{m,N}$ (rounded)\\
\hline
4  & 1  & 7     & 0.641409\\
4  & 3  & 27    & 15.837548\\
8  & 8  & 223   & 197.775136\\
16 & 16 & 1043  & 993.168330\\
32 & 32 & 5121  & 5022.804049\\
64 & 64 & 24705 & 24508.197207\\
\hline
\end{tabular}
\end{center}

Only the logarithms in the table are rounded.  For example, at
$(n,m)=(4,3)$ the primes $5$ and $7$ from the harmonic
denominator both disappear from the final reduced denominator.
The exact test still fails to produce a small linear form.
Thus the computation does account for actual cancellation.
Its finite range cannot decide the behavior of all $(m,N)$
or of combinations of different approximants.

\subsection{The exact lattice of one rank-three CMF trajectory}
\label{ga:cmf}

The family of Ap\'ery-like recurrences in \cite{CMF2026} supplies a
different setting in which denominator normalization matters.
The following calculation concerns the specific trajectory
$(0,0;1,1,1)$ and the projection used in \cite{VMRC033}.
It is a theorem about its complete image at each depth, not a claim
that every depth-dependent input inherits a fixed-input error bound.

Define the integer matrices
\[
 M_n=
 \begin{pmatrix}
 1&n^3&3(n+1)(n-1)(n^2+n+1)\\
 2&1-3n^2&-(n+1)(8n^2+n+5)\\
 1&3n+2&6n(n+1)
 \end{pmatrix}
\]
and
\[
 N_n=
 \begin{pmatrix}
 6n^2+13n+10&-8n-5&3\\
 3n^3+3n^2-3n-6&-3n^2+3n+3&n-2\\
 n^4-2n^3-6n^2-3n+3&6n^2+2n-1&1-2n
 \end{pmatrix}.
\]
Direct multiplication gives
\begin{equation}\label{ga:matrix-relations}
 \det M_n=\det N_n=(n+1)^6,\qquad
 N_n=(n+1)^4M_n^{-T}.
\end{equation}
Let $P_n=N_0\cdots N_{n-1}$, $U_n=M_0\cdots M_{n-1}$,
with both empty products equal to $I_3$, and put
\[
 R=\begin{pmatrix}1&1&0\\1&1&1\end{pmatrix},\qquad
 RP_n=\begin{pmatrix}\mathbf p_n\\\mathbf q_n\end{pmatrix},
 \qquad \Lambda_n=RP_n\Z^3.
\]

\begin{theorem}[Exact image and global content]
\label{ga:lattice}
With $H_0=0$, one has
\begin{equation}\label{ga:lattice-image}
 \Lambda_n=
 \{(n!a+3n!H_nb,\ n!b)^T:a,b\in\Z\}.
\end{equation}
If
$D_n=\operatorname{den}(3H_n)$ is the reduced positive denominator
and $C_n=3D_nH_n$, then the gcd of all entries of $RP_n$ is
$n!/D_n$.  Dividing the entire matrix by that common gcd gives
the image lattice
\begin{equation}\label{ga:normalized-lattice}
 \{(D_na+C_nb,\ D_nb)^T:a,b\in\Z\},
 \qquad [\Z^2:\text{this lattice}]=D_n^2.
\end{equation}
\end{theorem}
\begin{proof}
We give the two arithmetic components of the proof.
First, \eqref{ga:matrix-relations} implies
$P_n=(n!)^4U_n^{-T}$ and $\det P_n=(n!)^6$.
The cross product of the two rows of $R$ is $(1,-1,0)$, so
\begin{equation}\label{ga:cross}
 \operatorname{cross}(\mathbf p_n,\mathbf q_n)
       =(n!)^2\mathbf b_n,\qquad
 \mathbf b_n=(1,-1,0)U_n\in\Z^3.
\end{equation}
The vector $\mathbf b_n$ is primitive.  To see this, put
$\theta=z\,d/dz$, $\mathscr D=\theta-z$, and
\[
 J_n(z)=\left[\prod_{j=0}^{n-1}(j-\mathscr D)\right]^3 1.
\]
The gauge identity $e^z\theta e^{-z}=\mathscr D$ gives
\[
 e^{-z}J_n(z)
 =z^n\sum_{k\ge0}\frac{(-1)^k(n+k)!^2}{k!^3}z^k.
\]
This function satisfies
$((\theta-n)^3+z(\theta+1)^2)f=0$ and is carried to the
next one by $z(\theta+1)^2$.
Computing the first three jets at $z=1$ gives
\[
 \mathbf b_n=(J_n(1),\mathscr D J_n(1),\mathscr D^2J_n(1)).
\]
Over $\mathbb F_\ell[z]$ one has
$\mathscr D^\ell-\mathscr D=-z^\ell$.
For example, the left side commutes with multiplication by $z$,
and its value on $1$ is $-z^\ell$ by the Stirling-number formula
for $\mathscr D^\ell1$ and Fermat's identity.
It follows that
$\prod_{j=0}^{\ell-1}(j-\mathscr D)=z^\ell$.
Writing $n=a\ell+r$ with $0\le r<\ell$ therefore gives
$\mathbf b_n\equiv\mathbf b_r\pmod\ell$.
Every $M_j$, $0\le j<r$, is invertible modulo $\ell$, so
$\mathbf b_r\not\equiv0$.  No prime divides all entries of
$\mathbf b_n$.
By \eqref{ga:cross}, the gcd of the maximal minors of $RP_n$
is $(n!)^2$, which is the covolume of $\Lambda_n$.

Second, the harmonic shear of the two rows is integral.
Here are finite formulas making this assertion explicit.  Put
$D=u\,d/du$ and
\[
 Q_n(u)=\sum_{k=0}^n\binom nk^3 k!\,u^k,\qquad
 A_n(u)=\sum_{k=0}^n\binom nk^3 k!
          \{3(H_{n-k}-H_n)-2H_k\}u^k.
\]
Both polynomials have integer coefficients:
$k!H_k$ is integral, and
$\binom nk^3k!=\binom nk^2n!/(n-k)!$
is divisible by each integer between $n-k+1$ and $n$.
Set $a_n=1-3n$, $d_n=1+3n^2$, and define integer rows
\[
 \mathcal R_{Q,n}
 =\left((D^2+a_nD+d_n)Q_n,\ (D+a_n)Q_n,\ Q_n\right)_{u=1},
\]
\begin{equation}\label{ga:rowA}
\begin{split}
 \mathcal R_{A,n}={}&
 \bigl((D^2+a_nD+d_n)A_n+2DQ_n+a_nQ_n,\\
 &\hspace{15mm}(D+a_n)A_n+Q_n,\ A_n\bigr)_{u=1}.
\end{split}
\end{equation}
These rows satisfy
\begin{equation}\label{ga:row-covariance}
\begin{split}
 \mathcal R_{Q,n+1}&=\frac{\mathcal R_{Q,n}N_n}{n+1},\\
 \mathcal R_{A,n+1}&=\frac{\mathcal R_{A,n}N_n}{n+1}
                  -\frac{3\mathcal R_{Q,n+1}}{n+1}.
\end{split}
\end{equation}
One verification is to use the adjoint solutions
$w_Q=Q_n(1/z)$ and
$w_A=A_n(1/z)-\log z\,Q_n(1/z)$ of
$z\theta^2w-(\theta+n)^3w=0$.
For $f$ solving $((\theta-n)^3+z(\theta+1)^2)f=0$,
the expression
\[
 (\theta^2w-(z-3n)\theta w+(z+3n^2)w)f
 +(-\theta w+(z-3n)w)\theta f+w\theta^2f
\]
has derivative zero.  If
$f=z^n(\alpha(\log z)^2/2+\beta\log z+\chi+O(z\log^2z))$,
its values for $w_Q,w_A$ are respectively
$n!\alpha$ and $n!(\beta-5H_n\alpha)$.
The map $f\mapsto z(\theta+1)^2f$ changes
$(\alpha,\beta)$ to
$((n+1)^2\alpha,(n+1)^2\beta+2(n+1)\alpha)$.
Evaluating the conserved expression at $z=1$ proves
\eqref{ga:row-covariance}.
The displayed adjoint solutions themselves follow by differentiating
at $\varepsilon=0$ the finite sum
$z^{-\varepsilon}\sum_{k=0}^nc_k(\varepsilon)z^{-k}$, where
$c_0=1$ and
$c_{k+1}/c_k=(n-\varepsilon-k)^3/(\varepsilon+k+1)^2$;
its residual under $(\theta+n)^3-z\theta^2$ is
$-\varepsilon^2z^{1-\varepsilon}
 -\varepsilon^3c_n(\varepsilon)z^{-n-\varepsilon}$.

The initial rows are
$\mathcal R_{Q,0}=(1,1,1)$ and
$\mathcal R_{A,0}=(1,1,0)$.
Induction in \eqref{ga:row-covariance} gives
\[
 \mathbf q_n=n!\mathcal R_{Q,n},\qquad
 \mathbf p_n-3H_n\mathbf q_n=n!\mathcal R_{A,n}.
\]
Hence $\Lambda_n$ is contained in the right side of
\eqref{ga:lattice-image}.  Both lattices have covolume $(n!)^2$,
so they are equal.  Finally
$\gcd(n!,3n!H_n)=n!/D_n$ proves the content and normalization
assertions.
\end{proof}

\begin{remark}[Global normalization versus an individual reduced pair]
\label{ga:individual}
The division in Theorem~\ref{ga:lattice} removes the gcd shared by
\emph{all} entries of the projected matrix.  A selected pair
$(D_na+C_nb,D_nb)$ can have an additional gcd.
For example, if $\gcd(a,b)=1$ and this additional gcd is $g$, then
$g\mid D_n^2$: multiplication by the adjugate of
$\left(\begin{smallmatrix}D_n&C_n\\0&D_n\end{smallmatrix}\right)$
shows that $g$ divides both $D_n^2a$ and $D_n^2b$.
It would be incorrect to call every globally normalized output
already primitive.  Any method using additional per-output
cancellation must analyze that cancellation together with its error.
\end{remark}

At each fixed $n$, the nonzero-denominator ratios in
\eqref{ga:normalized-lattice} are exactly
$3H_n+a/b$, hence run through all of $\Q$.
Unrestricted freedom to change the input is therefore not, on its own,
an arithmetic approximation theorem.
The globally normalized linear form is
\begin{equation}\label{ga:CMF-form}
 \ell_n(a,b)=D_na+C_nb-\gamma D_nb
           =D_n\{a-b(\gamma-3H_n)\}.
\end{equation}
Dirichlet's theorem provides $a,b\in\Z$, $1\le b\le B$,
with $|\ell_n(a,b)|\le D_n/B$.
This upper bound tends to zero if $B/D_n\to\infty$.
On the other hand, under a reductio assumption $\gamma=r/s$,
$\ell_n(a,b)=0$ would imply $D_n\mid sb$, since
$\gcd(C_n,D_n)=1$.
Thus $0<|b|<D_n/s$ is a sufficient certificate of nonvanishing.
Using $B<D_n/s$ as a uniform certificate for the Dirichlet choices
is incompatible with making the bound $D_n/B$ tend to zero.

This last observation concerns these two sufficient estimates.
It does not rule out stronger approximation estimates, separately
proved nonvanishing, useful per-output gcd cancellation, or carefully
restricted adaptive inputs.  Under $\gamma=r/s$, every nonzero
integer-coefficient form is at least $1/s$ in absolute value,
while the exact lattice also contains nonzero pairs with form zero.
An argument that produces a small form but permits that form to
vanish has not supplied Proposition~\ref{ga:irrationality}.

\section{The remaining transfer requirements}
\label{end:interfaces}

The preceding sections provide several complete statements about objects
surrounding $\gamma$ and $M$.  Their implications for the two original
targets require additional inputs that must not be supplied by analogy.
This section records those inputs explicitly.

\subsection{Euler's constant}

The field-theoretic radical criterion would prove transcendence of
$\gamma$ if one could establish $\delta^{1/n}\in\K$ for some $n\ge2$.
No such membership is proved or predicted here.  The relative-closure
theorem controls elements already known to lie in $\K$; it does not force
an arbitrary analytic constant or one of its roots to belong to $\K$.
Likewise, $\delta=e(\eta-\gamma)$ shows that algebraic $\gamma$ would put
$\delta$ in the $E$-value ring, but that membership alone does not violate
any of the proved closure statements.

The reciprocal-series route requires the actual identity in
Corollary~\ref{rs:gamma-link}, with integer coefficients, the stated
size bound, and infinitely many nonzero coefficients.  The irrationality
of the permitted class of series does not identify $\gamma$ as a member.

For rational approximants the complete success criterion remains
\[
 P_j\in\Z,\quad Q_j\in\N,\quad
 0<|Q_j\gamma-P_j|\longrightarrow0.
\]
The Pad\'e family supplies logarithm cancellation and nonvanishing for
infinitely many indices.  The slow-degree theorem rules out the displayed
asymptotic degree range, and the 384-case certificate checks a finite
collection of larger degrees.  Neither result decides every remaining
degree schedule or every linear combination of approximants.  In the
CMF discussion, additional cancellation in a selected output pair has to
be analyzed together with its error and its nonvanishing; global matrix
normalization is not that analysis.

\subsection{The least classical Mills constant}

Under hypothetical algebraicity, the least-prime property and the cubic
trace formulas force special prime-free intervals immediately above
prime cubes.  An unconditional prime in one such interval, far enough
along the hypothetical tower, would contradict minimality.  The given
gap divisibility, exponent, three-adic, and splitting restrictions have
not been shown mutually inconsistent.

The shifted clocks supply fixed-polynomial escape, but applying it to a
classical algebraic candidate requires preserving the appropriate prime
labels at a scale where that escape estimate is effective.  The estimates
audited here do not furnish such a scale.  Taking limits of shifted
constants or discarding intermediate prime constraints does not repair
that missing step.

\subsection{Reusable completed statements}

The relative-closure and intersection results, the explicit compositum
counterexample, the reciprocal-series linear independence theorem, the
bounded-coefficient irrationality criterion, and the exact cubic trace
identities stand on their own stated hypotheses.  The accompanying
certificates establish the specified finite claims with integer or
rational arithmetic.  These are the completed results of this
compilation.  The two original arithmetic questions remain outside those
conclusions.

\appendix
\section{Reproducing the Pad\'e arithmetic certificate}
\label{app:pade}

The following Python program uses only the standard library.  Save the
listing as \texttt{pade\_certificate.py} and run
\texttt{python3 pade\_certificate.py 64} with Python 3.11 or later.
It recomputes all Bernoulli moments and all rational approximants; it
does not read precomputed numerical values of $\gamma$ or a separate
input dataset.  Every inequality is tested on exact integers through
\texttt{fractions.Fraction}.  The floating-point logarithms are only
display values.  An assertion checks that every tested lower bound is
greater than one.  The JSON output is generated by the program and is
not required to compile or read this manuscript.

The complete 384-case program was rerun for this edition and agreed
entry-for-entry with the preceding computation.  Its summary is
\[
 \text{cases}=384,\qquad
 \text{all }L_{m,N}>1,\qquad
 \min\log_{10}L_{m,N}=0.641409\ \text{(rounded)}.
\]
The unique source file contains the complete listing, not a reference to
an external script.


\begin{lstlisting}[style=certificate]

from fractions import Fraction as F
from math import comb, gcd, log10, isqrt
import json
import sys
import time
sys.set_int_max_str_digits(0)

max_m = int(sys.argv[1]) if len(sys.argv)>1 else 64
scales = [2,4,8,16,32,64]

def bernoullis(n):
    out = [F(1)]
    for j in range(1,n+1):
        out.append(-sum((F(comb(j+1,k))*out[k] for k in range(j)),F(0))/F(j+1))
    return out

def primes(a,b):
    return [p for p in range(a,b+1) if p>=2 and all(p%d for d in range(2,isqrt(p)+1))]

def harmonic(n):
    return sum((F(1,k) for k in range(1,n+1)),F(0))

def logf(q):
    return log10(q.numerator)-log10(q.denominator)

started=time.monotonic()
B=bernoullis(4*max_m+4)
mu=[abs(B[2*j+2])/F(2*j+2) for j in range(2*max_m+2)]
p_prev=[F(0)]
p=[F(1)]
h_prev=None
aa=[]
bb=[F(0)]
R={n:[] for n in scales}
prev={s:(F(0),F(1),F(1),F(0)) for n in scales for s in (n*n,4*n*n)}
history={s:[] for s in prev}
output=[]
for m in range(max_m+1):
    h=sum((p[i]*mu[i+m] for i in range(m+1)),F(0))
    a=sum((p[i]*mu[i+m+1] for i in range(m+1)),F(0))/h+(p[m-1] if m else 0)
    assert h>0 and a>0
    if m:
        b=h/h_prev
        assert b>0
        bb.append(b)
    else:
        b=F(0)
    aa.append(a)
    for s,(N1,N0,D1,D0) in list(prev.items()):
        # m=0 constructs R_1; m>=1 advances the usual Jacobi convergents.
        if m==0:
            N2,D2=mu[0],F(s)+a
        else:
            N2=(F(s)+a)*N1-b*N0
            D2=(F(s)+a)*D1-b*D0
        r=N2/D2
        if history[s]: assert r>history[s][-1]
        history[s].append(r)
        prev[s]=(N2,N1,D2,D1)
    # After R_{m+1}, compare the previous R_m at each scale.
    if m>=1:
        for n in scales:
            N=n.bit_length()-1
            rm=(N+1)*(harmonic(n)-F(1,2*n)+history[n*n][m-1])-N*(harmonic(2*n)-F(1,4*n)+history[4*n*n][m-1])
            Q=rm.denominator
            diff=history[n*n][m]-history[n*n][m-1]
            exact_lower=Q*diff
            interval_primes=primes(n+1,2*n)
            removed=[p for p in interval_primes if Q%p]
            row={'n':n,'m':m,'Q_digits':len(str(Q)),
                 'log10_certified_lower':round(logf(exact_lower),6),
                 'certified_lower_gt_1':exact_lower>1,
                 'fresh_primes':interval_primes,'removed_fresh_primes':removed}
            output.append(row)
            if m in {max(1,n//4),max(1,n//2),n,2*n,4*n}:
                print(json.dumps(row),flush=True)
    p_new=[F(0)]*(m+2)
    for i,co in enumerate(p):
        p_new[i+1]+=co
        p_new[i]-=a*co
    if m:
        for i,co in enumerate(p_prev): p_new[i]-=b*co
    p_prev,p=p,p_new
    h_prev=h
assert len(output) == max_m * len(scales)
assert all(x['certified_lower_gt_1'] for x in output)
print(json.dumps({'computed_max_m':max_m,'seconds':round(time.monotonic()-started,3),
                  'cases':len(output),'all_certified_lower_gt_1':all(x['certified_lower_gt_1'] for x in output),
                  'smallest_log10_lower':min(x['log10_certified_lower'] for x in output)}),flush=True)
with open('pade_certificate_results.json','w') as file:
    json.dump(output,file,indent=2)

\end{lstlisting}


\clearpage
\section{Reproducing the prime-seed nonreturn certificate}
\label{mi:app:nonreturn}

Save the following standard-library program as
\texttt{mills\_nonreturn\_verify.py} and run it with Python 3.11 or later.
All root brackets, primality checks, congruences, and orbit statements
are verified exactly. The loop checks every one of the $85\,712$
distinct states, not a sample of the cycle. The certificate establishes
nonreturn modulo the specified prime; it does not assert that all
subsequent traces are prime. In fact, the last assertions give the
explicit factorization of $T_2$.

\begin{lstlisting}[style=certificate]

"""Exact certificate for the nonreturning cubic prime seed.

Python 3 standard library only.  No floating-point arithmetic and
no probabilistic primality tests are used.
"""
from fractions import Fraction
from math import isqrt


def f(x):
    return x**3 - 211*x**2 - 7*x + 1


def is_prime(n):
    if n < 2:
        return False
    if n % 2 == 0:
        return n == 2
    return all(n % d for d in range(3, isqrt(n) + 1, 2))


def step(t, b, modulus=None):
    out = (t**3 - 3*t*b - 3, b**3 + 3*t*b + 3)
    if modulus is None:
        return out
    return tuple(x % modulus for x in out)


assert f(Fraction(-1, 10)) < 0 < f(0)
assert f(0) > 0 > f(1)
assert f(211) < 0 < f(212)
assert f(1) != 0 and f(-1) != 0
assert 212**17 < 10**40

t0, b0 = 211, -7
p, b1 = step(t0, b0)
assert (p, b1) == (9398359, -4771)
assert is_prime(t0) and is_prime(p)

roots = (-1584786, -1942624, 3527621)
assert len({x % p for x in roots}) == 3
assert all(f(x) % p == 0 for x in roots)

start = step(0, b1 % p, p)
assert start == (9398356, 7830596)
state = start
seen = set()
for _ in range(85712):
    assert state not in seen
    seen.add(state)
    assert state[0] != 0
    state = step(*state, p)
assert state == start
assert len(seen) == 85712

t2, _ = step(p, b1)
assert t2 == 830149079789343894643
assert t2 == 7 * 31 * 3825571796264257579
assert (211 - 1) % 3 == 0
assert (-7 + 1) % 3 == 0
assert (211 - 7) % 3 == 0 and (211 - 7) % 9 != 0

print("PASS: exact root brackets, primality, complete splitting,")
print("85712 distinct nonzero-trace cycle states, and T_2 factorization.")

\end{lstlisting}

\clearpage
\begin{thebibliography}{99}
\phantomsection
\addcontentsline{toc}{section}{References}

\bibitem{Beukers2006}
F.~Beukers, \emph{A refined version of the Siegel--Shidlovskii theorem},
Annals of Mathematics (2) \textbf{163} (2006), 369--379.
\url{https://doi.org/10.4007/annals.2006.163.369}.

\bibitem{FR2026}
S.~Fischler and T.~Rivoal, \emph{The ring of values of $E$-functions is not
a field}, preprint, October 2026, arXiv:2610.16603.
\url{https://arxiv.org/abs/2610.16603}.
Author version: \url{https://www.imo.universite-paris-saclay.fr/~fischler/Ecroix.pdf}.

\bibitem{NISTDLMF}
NIST Digital Library of Mathematical Functions, Sections 5.9 and 6.2,
especially equations 5.9.15, 5.9.16, and 6.2.4.
\url{https://dlmf.nist.gov/5.9}; \url{https://dlmf.nist.gov/6.2}.
Accessed 5 November 2026.

\bibitem{Saito2025}
K.~Saito, \emph{Mills' constant is irrational}, Mathematika
\textbf{71} (2025), no.~3, e70027.
\url{https://doi.org/10.1112/mtk.70027};
\url{https://arxiv.org/abs/2404.19461}.

\bibitem{SaitoVariants}
K.~Saito, \emph{Transcendency of variants of Mills' constant},
arXiv:2508.16068v3, 7 December 2025, especially Theorems 1.7 and 1.8.
\url{https://arxiv.org/abs/2508.16068}.

\bibitem{BHP2001}
R.~C.~Baker, G.~Harman, and J.~Pintz,
\emph{The difference between consecutive primes, II},
Proceedings of the London Mathematical Society (3) \textbf{83}
(2001), 532--562.
\url{https://doi.org/10.1112/plms/83.3.532}.

\bibitem{Schneider1934}
T.~Schneider, \emph{Transzendenzuntersuchungen periodischer Funktionen I.
Transzendenz von Potenzen}, Journal f\"ur die reine und angewandte
Mathematik \textbf{172} (1934), 65--69.
\url{https://eudml.org/doc/149900}.

\bibitem{CMF2026}
R.~Kalisch, T.~Monderer, E.~Leibtag, S.~Weinbaum, M.~Shalyt, and
I.~Kaminer, \emph{An Infinite Family of Ap\'ery-like Recurrences for
Euler's Constant}, arXiv:2610.13919v1, 9 October 2026.
\url{https://arxiv.org/abs/2610.13919}.
The paper explicitly does not claim an irrationality proof for $\gamma$.

\bibitem{OpenAIQRH}
OpenAI, \emph{The quasi-Riemann hypothesis: scope of the Lean formalization},
\texttt{openai/math}, \texttt{lean/docs/003.md}, accessed 5 November 2026.
\url{https://github.com/openai/math/blob/main/lean/docs/003.md}.
Used only to identify the stated $7/8$ scope, not as a substitute for DH.

\bibitem{VMRC008}
Verifiable Mathematical Reasoning Corpus, manuscript [008],
\emph{Additive Clocks, Non-Reversible Recurrences, and Transcendental
Prime-Representing Constants}, working manuscript, snapshot consulted
5 November 2026.  Source input for shifted constants.

\bibitem{VMRC010}
Verifiable Mathematical Reasoning Corpus, manuscript [010],
\emph{Arithmetic Constraints on a Hypothetically Algebraic Mills Constant:
Gap Divisibility, Transcendental Exponents, and Cubic Structures}, working
manuscript, snapshot consulted 5 November 2026.

\bibitem{VMRC019}
Verifiable Mathematical Reasoning Corpus, manuscript [019],
\emph{Euler's Constant: Complete Integrated Research Manuscript
(Exponential-Integral Foundations, Positive Moments, Infinite Families,
Local Arithmetic, and Transfer Obstructions)}, working manuscript,
snapshot consulted 5 November 2026.

\bibitem{VMRC030}
Verifiable Mathematical Reasoning Corpus, manuscript [030],
\emph{Arithmetic and Observation Classifications, Recovery, and
Transcendence Barriers}, working manuscript, snapshot consulted
5 November 2026.  Source for finite-phase descent and closure arguments.

\bibitem{VMRC031}
Verifiable Mathematical Reasoning Corpus, manuscript [031],
\emph{Prime-Representing Constants and Related Structures: Encodings,
Arithmetic, Geometry, and Modular Obstructions}, working manuscript,
snapshot consulted 5 November 2026.

\bibitem{VMRC033}
Verifiable Mathematical Reasoning Corpus, manuscript [033],
\emph{Euler's Constant: Arithmetic, Monodromy, and Gamma-Related
Structures}, working manuscript, snapshot consulted 5 November 2026.

\bibitem{VMRC049}
Verifiable Mathematical Reasoning Corpus, manuscript [049],
\emph{Winding Embeddings and Prime-Power Pisot Floors: Integral
Intersections, Composite Windows, and Local Obstructions}, working
manuscript, snapshot consulted 5 November 2026.

\bibitem{VMRC064}
Verifiable Mathematical Reasoning Corpus, manuscript [064],
\emph{Ternary Trace Dynamics and Shifted Mills Constants}, working
manuscript, snapshot consulted 5 November 2026.  Source input for the
$3^k-1$ and $3^k+1$ transcendence assertions.

\bibitem{VMRC075}
Verifiable Mathematical Reasoning Corpus, manuscript [075],
\emph{Shifted Mills and Cubic Pisot Traces}, working manuscript,
snapshot consulted 5 November 2026.

\bibitem{VMRC}
Verifiable Mathematical Reasoning Corpus, repository for the working
manuscripts cited above.
\url{https://github.com/QEEDDD/Verifiable-Mathematical-Reasoning-Corpus}.
The corpus citations record the provenance of internal arguments; they
are not claims of external peer review or independent historical priority.

\end{thebibliography}
\end{document}
